\providecommand{\SO}{\operatorname{SO}}
\providecommand{\OO}{\operatorname{O}}
\providecommand{\SE}{\operatorname{SE}}
\providecommand{\Eucl}{\operatorname{E}}
\providecommand{\GL}{\operatorname{GL}}
\providecommand{\Sym}{\mathfrak{S}}
\providecommand{\Haar}{\mathrm{d}\mu}
\providecommand{\softmax}{\operatorname{softmax}}
\providecommand{\Stab}{\operatorname{Stab}}
\providecommand{\Orb}{\operatorname{Orb}}
\providecommand{\Attn}{\operatorname{Attn}}
\providecommand{\MLP}{\operatorname{MLP}}
\providecommand{\diag}{\operatorname{diag}}

\chapter{Groups: Group-Equivariant Networks}
\label{ch:groups}

The convolutional networks of \Cref{ch:grids} owe their success to a single
symmetry: translation. A filter that recognizes an edge in one corner of an
image recognizes the same edge anywhere, because convolution commutes with
shifts. Yet translation is only one symmetry among many. A microscope image of
tissue carries no privileged ``up''; a satellite photograph may be rotated by
any angle; a molecule may be reflected into its mirror image. A network that is
equivariant only to translations must rediscover each rotated or reflected
pattern independently, wasting capacity on copies of the same feature.

This chapter develops the systematic answer: replace the translation group by an
arbitrary symmetry group $\group$ and generalize convolution accordingly. The
result --- the \emph{group convolution} --- yields networks that are
equivariant by construction to rotations, reflections, and other geometric
transformations. We build the theory from finite groups such as the symmetries
of the square, extend it to continuous compact groups ($\SO(2)$, $\SO(3)$) via
steerable filters and harmonic analysis, push it to the non-compact Euclidean
groups ($\Eucl(3)$, $\SE(3)$) that govern three-dimensional point clouds and
molecules, and close with the permutation group, the symmetry of unordered
sets. Throughout, the guiding principle established in \Cref{ch:priors} remains
in force: equivariant layers compose into equivariant networks, and a single
invariant readout at the end produces a prediction insensitive to the chosen
pose.

\section{Beyond translations}
\label{sec:groups:beyond}

Recall from \Cref{ch:priors} that a map $\Phi$ between spaces of signals is
\emph{$\group$-equivariant} if $\Phi(L_g f) = L'_g\,\Phi(f)$ for every
$g \in \group$, where $L_g$ and $L'_g$ denote the action of $g$ on the input and
output spaces, and \emph{$\group$-invariant} if $\Phi(L_g f) = \Phi(f)$.
Equivariance generalizes the translation invariance of classical convolutional
networks: if the data carry symmetries under $\group$, then $\group$-equivariant
maps preserve those symmetries throughout processing, yielding representations
that are both more robust and more parameter-efficient. The mathematical
language of groups, group actions, and representations is developed in
\Cref{ch:foundations}; here we put it to work to construct architectures.

The classical convolutional network is the special case $\group = (\Z^2, +)$
acting by translations on discrete images. In this case the irreducible
representations are one-dimensional (the characters $x \mapsto e^{2\pi i
\langle k, x\rangle}$), group convolution reduces to ordinary convolution, and
filters are translation-invariant. The geometric extensions of this chapter
incorporate richer groups: $\SO(2)$ for planar rotations, $\SE(2)$ for rigid
planar motions, $\SO(3)$ and $\SE(3)$ for three-dimensional rotations and rigid
motions, and the symmetric group $\Sym_n$ for permutations of unordered
collections. The design recipe is always the same:
\begin{enumerate}
  \item \textbf{Identify the group}: determine which symmetries are present in
        the data.
  \item \textbf{Analyze the representations}: characterize the relevant
        irreducible representations under which features transform.
  \item \textbf{Design the convolution}: implement equivariant operations
        efficiently using the group structure.
  \item \textbf{Optimize}: exploit the symmetry to share parameters and reduce
        the hypothesis space without loss of expressive power.
\end{enumerate}

\section{Group convolutions}
\label{sec:groups:convolutions}

\subsection{Convolution on a group}
\label{sec:groups:conv-on-group}

The generalization of convolution to an arbitrary group is the theoretical core
of the geometric extensions of convolutional networks. The foundations of group
representation theory underlying this construction trace back to the classical
work of \citet{Peter1927} and \citet{Serre1977}.

\begin{definition}[Group convolution]
\label{def:groups:group-convolution}
For functions $f, \psi: \group \to \R$ defined on a group $\group$, the
\emph{group convolution} is
\begin{equation}
  (f *_{\group} \psi)(u)
  = \int_{\group} f(uv^{-1})\,\psi(v)\,\Haar(v),
  \label{eq:groups:group-convolution}
\end{equation}
where $\mu$ is the Haar measure on $\group$.
\end{definition}

This definition generalizes ordinary convolution by replacing the subtraction
$u - v$ with the group operation $uv^{-1}$. For discrete groups the integral
becomes a sum.

The following theorem establishes the deep connection between equivariance and
convolution: convolution is not merely \emph{one} way to build an equivariant
layer, it is essentially the \emph{only} way.

\begin{theorem}[Kondor--Trivedi]
\label{thm:groups:kondor-trivedi}
A feedforward neural network $N$ is equivariant to the action of a compact group
$\group$ on its inputs if and only if each layer of $N$ implements a generalized
form of convolution derived from \Cref{eq:groups:group-convolution}
\citep{kondor2018generalizationconvolutionalneural}.
\end{theorem}

\begin{proof}[Sketch]
The proof has two directions. \emph{Forward:} if each layer implements a group
convolution, then by construction it preserves the group structure, which
implies equivariance (verified directly below in
\Cref{thm:groups:lift-equivariance,thm:groups:group-conv-equivariance}).
\emph{Converse:} if the network is equivariant, then representation theory
characterizes all admissible linear equivariant maps as generalized
convolutions. The full argument uses harmonic analysis on $\group$ and the
Peter--Weyl decomposition to describe every equivariant linear map between
feature spaces.
\end{proof}

The choice of representation under which features transform determines what
kinds of patterns a layer can detect equivariantly. In a $\group$-equivariant
network, each layer processes features that transform under a specific
representation $\rho: \group \to \GL(V)$ of the symmetry group.

\begin{example}[Regular representation in group convolutional networks]
\label{ex:groups:regular-representation}
For a finite group $\group = \{g_1, \dots, g_n\}$, the \emph{left regular
representation} acts on functions $f: \group \to \R^d$ by
\begin{equation}
  (\rho(g) f)(h) = f(g^{-1} h).
\end{equation}
This representation is fundamental in group convolutional networks, where the
channels correspond to evaluations of the function at the different elements of
the group.
\end{example}

By the Peter--Weyl theorem (\Cref{ch:foundations}), the space $L^2(\group)$ of a
compact group decomposes as $L^2(\group) = \bigoplus_{\lambda \in \widehat
\group} n_\lambda \mathcal{H}_\lambda$, a direct sum over the irreducible
representations $\mathcal{H}_\lambda$ with multiplicities $n_\lambda$. This has a
direct consequence for architecture design: each $\mathcal{H}_\lambda$
corresponds to a channel type, and $n_\lambda$ prescribes how many channels of
that type to include in order to capture the group symmetries in full.

In practice, signals are not defined on the whole group but on a
\emph{homogeneous space} $X = \group / H$, where $H$ is a subgroup.

\begin{definition}[Homogeneous space]
\label{def:groups:homogeneous-space}
A space $X$ is \emph{homogeneous} under the action of $\group$ if for any
$x, y \in X$ there exists $g \in \group$ with $g \cdot x = y$.
\end{definition}

The sphere $S^2$ is homogeneous under $\SO(3)$, and ordinary images are
homogeneous under translations. Generalizing convolution to a homogeneous space
requires additional \emph{lifting} and \emph{projection} machinery, developed
next.

\subsection{Lifting and group convolution layers}
\label{sec:groups:lifting}

A group convolutional network proceeds in three stages: a \emph{lifting layer}
maps a planar signal to a function on the group, a stack of \emph{group
convolution layers} processes that function equivariantly, and an
\emph{invariant readout} collapses the group dimension into a prediction.

\paragraph{Lifting.}
The first operation transforms a signal on Euclidean space $\R^2$ into a
function on the group $\group$.

\begin{definition}[Lifting convolution]
\label{def:groups:lifting-convolution}
Given a function $f: \R^2 \to \R$ and a filter $\psi: \R^2 \to \R$, the
\emph{lifting convolution} produces a function $F: \group \to \R$,
\begin{equation}
  (f \star \psi)(g) = \int_{\R^2} f(x)\,\psi(g^{-1} \cdot x)\,dx,
  \label{eq:groups:lifting-conv}
\end{equation}
where $g \in \group$ and $g^{-1} \cdot x$ denotes the action of $g^{-1}$ on the
point $x \in \R^2$.
\end{definition}

For a finite group $\group$ the integral is a sum over a local window, so that
for each spatial position and each transformation $g \in \group$ we compute the
response of the filter $\psi$ transformed by $g^{-1}$.

\begin{theorem}[Equivariance of lifting]
\label{thm:groups:lift-equivariance}
The lifting convolution is equivariant in the sense that
\begin{equation}
  (L_h f \star \psi)(g) = (f \star \psi)(h^{-1} g),
\end{equation}
where $L_h f(x) = f(h^{-1} \cdot x)$ is the action of $\group$ on the input.
\end{theorem}

\begin{proof}
By definition,
\begin{align}
  (L_h f \star \psi)(g)
  &= \int_{\R^2} (L_h f)(x)\,\psi(g^{-1} \cdot x)\,dx
   = \int_{\R^2} f(h^{-1} \cdot x)\,\psi(g^{-1} \cdot x)\,dx.
\end{align}
With the change of variables $y = h^{-1} \cdot x$ (whose Jacobian is constant
because the action is by isometries),
\begin{align}
  &= \int_{\R^2} f(y)\,\psi(g^{-1} \cdot h \cdot y)\,dy
   = \int_{\R^2} f(y)\,\psi\big((h^{-1} g)^{-1} \cdot y\big)\,dy
   = (f \star \psi)(h^{-1} g). \qedhere
\end{align}
\end{proof}

\paragraph{Group convolution.}
Subsequent layers operate between functions on the group.

\begin{definition}[Group convolution layer]
\label{def:groups:group-conv-layer}
For functions $f, \psi: \group \to \R$, the \emph{group convolution} is
\begin{equation}
  (f \star_{\group} \psi)(u) = \sum_{g \in \group} f(g)\,\psi(g^{-1} u),
  \label{eq:groups:group-conv-discrete}
\end{equation}
where the sum ranges over all elements of $\group$. This is the discrete
counterpart of \Cref{eq:groups:group-convolution}: the change of variables
$v \mapsto g = uv^{-1}$ turns $\int_{\group} f(uv^{-1})\psi(v)\,\Haar(v)$ into
$\sum_g f(g)\psi(g^{-1}u)$.
\end{definition}

\begin{theorem}[Equivariance of group convolution]
\label{thm:groups:group-conv-equivariance}
Group convolution is equivariant:
\begin{equation}
  (L_h f) \star_{\group} \psi = L_h (f \star_{\group} \psi),
\end{equation}
where $L_h$ is the left-translation action $(L_h f)(g) = f(h^{-1} g)$.
\end{theorem}

\begin{proof}
By definition,
\begin{align}
  \big((L_h f) \star_{\group} \psi\big)(u)
  &= \sum_{g \in \group} (L_h f)(g)\,\psi(g^{-1} u)
   = \sum_{g \in \group} f(h^{-1} g)\,\psi(g^{-1} u).
\end{align}
With the substitution $g' = h^{-1} g$ (equivalently $g = h g'$),
\begin{align}
  &= \sum_{g' \in \group} f(g')\,\psi\big((h g')^{-1} u\big)
   = \sum_{g' \in \group} f(g')\,\psi\big(g'^{-1} h^{-1} u\big) \\
  &= (f \star_{\group} \psi)(h^{-1} u)
   = \big(L_h (f \star_{\group} \psi)\big)(u). \qedhere
\end{align}
\end{proof}

Group convolution inherits the algebra of the group algebra $\R[\group]$, in
which multiplication is convolution; this algebraic structure is what guarantees
that composing layers preserves equivariance.

\begin{algorithm}[h]
\caption{Group convolution layer for a finite group}
\label{alg:groups:gconv-finite}
\begin{algorithmic}[1]
\Require Input function $f$, filters $\{\psi_i\}$, group
  $\group = \{g_1, \dots, g_n\}$
\Ensure Feature map $\mathrm{output}$
\State Initialize $\mathrm{output}$ with shape $(|\group|, H, W, C_{\text{out}})$
\For{each filter $i \in \{1, \dots, C_{\text{out}}\}$}
  \For{each group element $u \in \group$}
    \For{each position $(h, w)$}
      \State $s \leftarrow 0$
      \For{each group element $g \in \group$}
        \For{each input channel $c$}
          \State $s \leftarrow s + f[g][h][w][c] \cdot \psi_i[u^{-1} g][c]$
        \EndFor
      \EndFor
      \State $\mathrm{output}[u][h][w][i] \leftarrow s$
    \EndFor
  \EndFor
\EndFor
\State \Return $\mathrm{output}$
\end{algorithmic}
\end{algorithm}

\paragraph{A worked example: the dihedral group $D_4$.}
The dihedral group $D_4$ of symmetries of the square is the paradigmatic finite
group for image data.

\begin{definition}[The group $D_4$]
\label{def:groups:d4}
The dihedral group $D_4$ has eight elements,
\begin{equation}
  D_4 = \{e, r, r^2, r^3, s, sr, sr^2, sr^3\},
\end{equation}
where $e$ is the identity, $r$ is a $90^\circ$ rotation ($r^4 = e$), $s$ is a
reflection ($s^2 = e$), and the relation $s r s = r^{-1}$ holds.
\end{definition}

\begin{proposition}[Action of $D_4$ on the plane]
\label{prop:groups:d4-action}
The natural action of $D_4$ on $\R^2$ is generated by
\begin{align}
  r \cdot (x,y) &= (-y, x) && \text{($90^\circ$ rotation),} \\
  s \cdot (x,y) &= (x, -y) && \text{(horizontal reflection).}
\end{align}
All other elements arise by composition.
\end{proposition}

\begin{proof}
The rotation $r$ acts via $R_{90} = \left(\begin{smallmatrix} 0 & -1 \\ 1 & 0
\end{smallmatrix}\right)$, giving $R_{90}(x,y)^\top = (-y,x)^\top$. The
reflection $s$ across the $x$-axis acts via $S = \left(\begin{smallmatrix} 1 & 0
\\ 0 & -1 \end{smallmatrix}\right)$, giving $(x,-y)^\top$. The remaining six
elements follow by composition: $r^2 \cdot (x,y) = (-x,-y)$, $r^3 \cdot (x,y) =
(y,-x)$, $sr \cdot (x,y) = (-y,-x)$, $sr^2 \cdot (x,y) = (-x,y)$, $sr^3 \cdot
(x,y) = (y,x)$. These eight maps are exactly the linear isometries of $\R^2$
preserving the unit square.
\end{proof}

The eight elements satisfy the group axioms (closure, associativity, identity,
and inverses), which may be verified by constructing the multiplication table
of $D_4$. \Cref{fig:groups:d4-actions} shows the eight transformations acting on
an asymmetric patch.

\begin{figure}[h]
\centering
\begin{tikzpicture}[
    filled/.style={minimum size=0.55cm, inner sep=0pt, outer sep=0pt,
                   fill=gdlblue!40, draw=gdlblue!60},
    empty/.style={minimum size=0.55cm, inner sep=0pt, outer sep=0pt,
                  fill=gdlgray!15, draw=gdlgray!40},
    grouplabel/.style={font=\large, text=gdlpurple!70!black},
    rowlabel/.style={font=\small\itshape, text=gdlgray!80!black, rotate=90,
                     anchor=south}
]
\def\cs{0.55}
\def\hs{3.0}
\def\vs{3.2}
% TOP ROW: rotations
\begin{scope}[shift={(0, 0)}]
    \node[filled] at (0,0){};\node[filled] at (\cs,0){};\node[filled] at (2*\cs,0){};
    \node[filled] at (0,-\cs){};\node[empty] at (\cs,-\cs){};\node[empty] at (2*\cs,-\cs){};
    \node[filled] at (0,-2*\cs){};\node[filled] at (\cs,-2*\cs){};\node[empty] at (2*\cs,-2*\cs){};
    \node[grouplabel] at (\cs, -2*\cs - 0.7) {$e$};
\end{scope}
\begin{scope}[shift={(\hs, 0)}]
    \node[filled] at (0,0){};\node[empty] at (\cs,0){};\node[empty] at (2*\cs,0){};
    \node[filled] at (0,-\cs){};\node[empty] at (\cs,-\cs){};\node[filled] at (2*\cs,-\cs){};
    \node[filled] at (0,-2*\cs){};\node[filled] at (\cs,-2*\cs){};\node[filled] at (2*\cs,-2*\cs){};
    \node[grouplabel] at (\cs, -2*\cs - 0.7) {$r$};
\end{scope}
\begin{scope}[shift={(2*\hs, 0)}]
    \node[empty] at (0,0){};\node[filled] at (\cs,0){};\node[filled] at (2*\cs,0){};
    \node[empty] at (0,-\cs){};\node[empty] at (\cs,-\cs){};\node[filled] at (2*\cs,-\cs){};
    \node[filled] at (0,-2*\cs){};\node[filled] at (\cs,-2*\cs){};\node[filled] at (2*\cs,-2*\cs){};
    \node[grouplabel] at (\cs, -2*\cs - 0.7) {$r^2$};
\end{scope}
\begin{scope}[shift={(3*\hs, 0)}]
    \node[filled] at (0,0){};\node[filled] at (\cs,0){};\node[filled] at (2*\cs,0){};
    \node[filled] at (0,-\cs){};\node[empty] at (\cs,-\cs){};\node[filled] at (2*\cs,-\cs){};
    \node[empty] at (0,-2*\cs){};\node[empty] at (\cs,-2*\cs){};\node[filled] at (2*\cs,-2*\cs){};
    \node[grouplabel] at (\cs, -2*\cs - 0.7) {$r^3$};
\end{scope}
\node[rowlabel] at (-1.1, -\cs) {Rotations};
% BOTTOM ROW: reflections
\begin{scope}[shift={(0, -\vs)}]
    \node[filled] at (0,0){};\node[filled] at (\cs,0){};\node[filled] at (2*\cs,0){};
    \node[empty] at (0,-\cs){};\node[empty] at (\cs,-\cs){};\node[filled] at (2*\cs,-\cs){};
    \node[empty] at (0,-2*\cs){};\node[filled] at (\cs,-2*\cs){};\node[filled] at (2*\cs,-2*\cs){};
    \node[grouplabel] at (\cs, -2*\cs - 0.7) {$s$};
\end{scope}
\begin{scope}[shift={(\hs, -\vs)}]
    \node[empty] at (0,0){};\node[empty] at (\cs,0){};\node[filled] at (2*\cs,0){};
    \node[filled] at (0,-\cs){};\node[empty] at (\cs,-\cs){};\node[filled] at (2*\cs,-\cs){};
    \node[filled] at (0,-2*\cs){};\node[filled] at (\cs,-2*\cs){};\node[filled] at (2*\cs,-2*\cs){};
    \node[grouplabel] at (\cs, -2*\cs - 0.7) {$sr$};
\end{scope}
\begin{scope}[shift={(2*\hs, -\vs)}]
    \node[filled] at (0,0){};\node[filled] at (\cs,0){};\node[empty] at (2*\cs,0){};
    \node[filled] at (0,-\cs){};\node[empty] at (\cs,-\cs){};\node[empty] at (2*\cs,-\cs){};
    \node[filled] at (0,-2*\cs){};\node[filled] at (\cs,-2*\cs){};\node[filled] at (2*\cs,-2*\cs){};
    \node[grouplabel] at (\cs, -2*\cs - 0.7) {$sr^2$};
\end{scope}
\begin{scope}[shift={(3*\hs, -\vs)}]
    \node[filled] at (0,0){};\node[filled] at (\cs,0){};\node[filled] at (2*\cs,0){};
    \node[filled] at (0,-\cs){};\node[empty] at (\cs,-\cs){};\node[filled] at (2*\cs,-\cs){};
    \node[filled] at (0,-2*\cs){};\node[empty] at (\cs,-2*\cs){};\node[empty] at (2*\cs,-2*\cs){};
    \node[grouplabel] at (\cs, -2*\cs - 0.7) {$sr^3$};
\end{scope}
\node[rowlabel] at (-1.1, -\vs - \cs) {Reflections};
\end{tikzpicture}
\caption[Elements of $D_4$ acting on a patch]{The eight elements of the dihedral
group $D_4$ acting on an asymmetric patch. Top row: the four rotations
($e$, $r$, $r^2$, $r^3$). Bottom row: the four reflections ($s$, $sr$, $sr^2$,
$sr^3$). A $D_4$-equivariant network produces the same output regardless of
which transformation is applied to the input.}
\label{fig:groups:d4-actions}
\end{figure}

A complete $D_4$-equivariant network composes a lifting layer, several group
convolution layers, and a final invariant layer.

\begin{definition}[$D_4$-equivariant network]
\label{def:groups:d4-network}
A network $\Phi: \mathcal{F}(\R^2) \to \R^C$ is \emph{$D_4$-equivariant} if it
decomposes as
\begin{equation}
  \Phi = \Phi_{\text{inv}} \circ \Phi_{\group} \circ \Phi_{\group}
         \circ \Phi_{\text{lift}},
\end{equation}
where $\Phi_{\text{lift}}: \mathcal{F}(\R^2) \to \mathcal{F}(D_4 \times \R^2)$ is
the lifting layer, $\Phi_{\group}: \mathcal{F}(D_4 \times \R^2) \to
\mathcal{F}(D_4 \times \R^2)$ are group convolution layers, and
$\Phi_{\text{inv}}: \mathcal{F}(D_4 \times \R^2) \to \R^C$ is a final invariant
layer.
\end{definition}

\begin{theorem}[Composition toward invariance]
\label{thm:groups:composition-invariance}
If each component $\Phi_i$ in \Cref{def:groups:d4-network} is $D_4$-equivariant
and the final layer $\Phi_{\text{inv}}$ is $D_4$-invariant, then the composition
$\Phi$ is $D_4$-invariant:
\begin{equation}
  \Phi(L_g f) = \Phi(f) \qquad \forall g \in D_4,\ f \in \mathcal{F}(\R^2).
\end{equation}
\end{theorem}

\begin{proof}
The composition of equivariant maps is equivariant (a property established in
\Cref{ch:priors}), so $\Phi_{\group} \circ \Phi_{\group} \circ \Phi_{\text{lift}}$
is equivariant. The final invariant layer then removes the group dependence,
producing an invariant output.
\end{proof}

In each equivariant layer the features are organized as functions $f: D_4 \times
\R^2 \to \R^d$; the group dimension ($|D_4| = 8$) introduces a multiplicative
factor in the feature space that preserves all information about orientations
and reflections.

\paragraph{Verifying equivariance.}
Theoretical correctness should be complemented by a numerical measure of how
well equivariance is preserved.

\begin{definition}[Equivariance error]
\label{def:groups:equivariance-error}
For a map $\Phi$ and a group $\group$, the \emph{equivariance error} at an input
$x$ is
\begin{equation}
  \mathcal{E}_{\group}(\Phi, x)
  = \frac{1}{|\group|} \sum_{g \in \group}
    \norm{\Phi(g \cdot x) - g \cdot \Phi(x)}^2,
\end{equation}
where $\norm{\cdot}$ is a suitable norm on the output space.
\end{definition}

\begin{theorem}[Equivariance criterion]
\label{thm:groups:equivariance-criterion}
A map $\Phi$ is $\group$-equivariant if and only if $\mathcal{E}_{\group}(\Phi,x)
= 0$ for every input $x$.
\end{theorem}

\begin{proof}
($\Rightarrow$) If $\Phi$ is equivariant then $\Phi(g \cdot x) = g \cdot
\Phi(x)$ for all $g$, so every summand vanishes and $\mathcal{E}_{\group}(\Phi,x)
= 0$. ($\Leftarrow$) If the average of the nonnegative terms is zero then each
term vanishes, so $\Phi(g \cdot x) = g \cdot \Phi(x)$ for all $g$ by positive
definiteness of the norm.
\end{proof}

\begin{proposition}[Error propagation]
\label{prop:groups:error-propagation}
For a composition $\Phi = \Phi_n \circ \cdots \circ \Phi_1$ in which each
$\Phi_i$ has Lipschitz constant $L_i$ and $x_i = (\Phi_{i-1} \circ \cdots \circ
\Phi_1)(x)$ are the intermediate activations (with $x_0 = x$), the worst-case
equivariance deviation propagates as
\begin{equation}
  \sup_{g \in \group} \norm{\Phi(g \cdot x) - g \cdot \Phi(x)}
  \leq \sum_{i=1}^{n} \Big(\prod_{j=i+1}^{n} L_j\Big)
    \sup_{g \in \group} \norm{\Phi_i(g \cdot x_i) - g \cdot \Phi_i(x_i)},
\end{equation}
where the empty product (for $i = n$) equals $1$.
\end{proposition}

\begin{proof}
For $n = 1$ the bound is trivial. For the inductive step let $\Psi =
\Phi_{n-1} \circ \cdots \circ \Phi_1$. Fixing $g \in \group$,
\begin{align}
  \norm{\Phi(g \cdot x) - g \cdot \Phi(x)}
  &= \norm{\Phi_n(\Psi(g \cdot x)) - g \cdot \Phi_n(\Psi(x))} \\
  &\leq \norm{\Phi_n(\Psi(g \cdot x)) - \Phi_n(g \cdot \Psi(x))}
      + \norm{\Phi_n(g \cdot \Psi(x)) - g \cdot \Phi_n(\Psi(x))}.
\end{align}
The first term is bounded by $L_n \norm{\Psi(g \cdot x) - g \cdot \Psi(x)}$
through Lipschitz continuity; the second is the equivariance deviation of
$\Phi_n$ at $x_n = \Psi(x)$. Taking the supremum over $g$ and applying the
inductive hypothesis to $\Psi$ yields the claim.
\end{proof}

\begin{remark}[Numerical stability of equivariance]
\label{rem:groups:numerical-stability}
For a $D_4$-equivariant architecture implemented in finite-precision arithmetic
the equivariance error is of order
\begin{equation}
  \sup_{g \in \group} \norm{\Phi(g \cdot x) - g \cdot \Phi(x)}
  = O\!\big(\epsilon_{\text{mach}} \cdot \operatorname{depth}(\Phi)
    \cdot \norm{x}\big),
\end{equation}
where $\epsilon_{\text{mach}}$ is the machine precision and
$\operatorname{depth}(\Phi)$ the network depth. Each layer introduces a rounding
error proportional to $\epsilon_{\text{mach}}$, and these accumulate linearly,
mirroring the standard error-propagation analysis of floating-point arithmetic.
Errors of order $10^{-6}$ observed in practice are consistent with 32-bit
floating-point precision: the theoretical equivariance is preserved up to the
hardware limit, and $\epsilon_{\text{mach}} \to 0$ recovers exact equivariance.
\end{remark}

\subsection{Why equivariance pays}
\label{sec:groups:why-pays}

A central result justifies the use of group convolutions: equivariant
architectures can approximate any continuous equivariant function, so symmetry
constrains the hypothesis space without sacrificing expressive power.

\begin{theorem}[Equivariant universality]
\label{thm:groups:universality}
Let $\group$ be a compact group acting on a compact domain $\mathcal{D}$, and let
$\mathcal{F}_{\group}$ be the space of continuous $\group$-equivariant maps
$L^2(\mathcal{D}) \to L^2(\mathcal{D})$. For every $f \in \mathcal{F}_{\group}$
and $\epsilon > 0$ there exists a $\group$-equivariant convolutional network
$\Phi_{\text{GCNN}}$ with $\norm{f - \Phi_{\text{GCNN}}}_{L^2} < \epsilon$. That
is, group convolutional networks are universal approximators in the space of
equivariant functions.
\end{theorem}

\begin{proof}[Sketch for $\group = \Z^2$]
For a compactly supported $f: \Z^2 \to \R$, the discrete Fourier transform gives
$f(x) = \sum_{k} \hat f(k)\, e^{2\pi i \langle k, x\rangle / N}$. Each frequency
component is approximated by a convolution with an appropriate filter, and a
convolutional network with $M$ first-layer filters can be written
$\Phi_{\text{CNN}}(f) = \sigma\big(\sum_{j=1}^M w_j (f * h_j) + b\big)$ for
learnable filters $h_j$, weights $w_j$, bias $b$, and activation $\sigma$.
Equivariance holds automatically because $L_g(f * h) = (L_g f) * h$, so
$\Phi_{\text{CNN}}(L_g f) = L_g \Phi_{\text{CNN}}(f)$. For non-abelian groups the
argument carries over by replacing the Fourier transform with harmonic analysis
on $\group$: decomposition into irreducible representations, group convolution,
and representation theory together guarantee that the network reproduces every
equivariant linear map.
\end{proof}

\begin{corollary}[Parameter efficiency]
\label{cor:groups:parameter-efficiency}
A $\group$-equivariant network requires $|\group|$ times fewer parameters than an
unconstrained network to approximate a $\group$-equivariant function to the same
accuracy, where $|\group|$ is the order of the group.
\end{corollary}

\begin{proof}
Let $f: X \to Y$ be $\group$-equivariant. An unconstrained network must learn the
value of $f$ at each point $x$ and at all its transforms $\{g \cdot x : g \in
\group\}$ independently, requiring $N$ parameters to reach accuracy $\epsilon$. A
$\group$-equivariant network satisfies $\Phi(g \cdot x) = g \cdot \Phi(x)$ by
construction, so it suffices to learn $\Phi$ on a fundamental domain of the
action, of size $|X|/|\group|$; the values on the remaining $|\group| - 1$ copies
follow by equivariance. This reduces the parameter count to $N/|\group|$ for the
same accuracy.
\end{proof}

This makes precise why equivariant architectures generalize better: they exploit
the symmetries of the data automatically, shrinking the hypothesis space without
loss of expressive power.

\paragraph{Finite groups in vision.}
The finite groups most relevant to image processing preserve the discrete pixel
grid while encoding meaningful symmetries: the cyclic group $C_4 = \{e, r, r^2,
r^3\}$ of $90^\circ$ rotations, the dihedral group $D_4$ adding reflections, and
the smaller groups $C_2$ (bilateral symmetry) and $D_2$. These groups preserve
the grid structure, appear naturally in textures, crystals, and medical images,
and admit exact implementations because they are finite. For the cyclic group
$C_4$, an equivariant network requires exactly one quarter of the parameters of
a standard network representing the same space of $C_4$-equivariant functions:
the four rotated copies $\{\psi, r\cdot\psi, r^2\cdot\psi, r^3\cdot\psi\}$ are
generated from a single base filter $\psi$ by the group action, a direct
instance of \Cref{cor:groups:parameter-efficiency}.

\begin{remark}[Dimensionality--expressivity trade-off]
\label{rem:groups:tradeoff}
For a group convolutional network processing functions on $\group \times
\Omega$, the feature-space dimension scales as $|\group| \cdot \dim(\Omega) \cdot
C$. The weight sharing induced by equivariance, however, reduces the number of
free parameters. If the data carry additional symmetries --- so the stabilizer
$\Stab(x) = \{g \in \group : g \cdot x = x\}$ is nontrivial for typical inputs
--- the effective parameter dimension is reduced further. By the
orbit--stabilizer theorem, $|\group| = |\Orb(x)| \cdot |\Stab(x)|$, which
suggests heuristically that the effective dimension is of order $|\Orb(x)| \cdot
\dim(\Omega) \cdot C$ rather than $|\group| \cdot \dim(\Omega) \cdot C$.
Computationally, the cost grows linearly with the group order, $\text{FLOPs}_{
\text{GCNN}} = |\group| \cdot \text{FLOPs}_{\text{CNN}}$, which can be
prohibitive for large groups.
\end{remark}

On rotated benchmarks the benefit is dramatic. The experiments of
\citet{cohen2016group} show that on a rotated version of CIFAR-10 a standard
convolutional network drops to about $47\%$ accuracy, whereas $C_4$- and
$D_4$-equivariant networks retain roughly $72\%$ and $73\%$, with a slight
improvement also on the unrotated task. The price is sensitivity to deviations
from exact symmetry: in real data the symmetry is often only approximate, so the
group must be chosen with care.

\section{Steerable filters and irreducible representations}
\label{sec:groups:steerable}

Finite groups capture discrete symmetries exactly, but a $C_4$-equivariant
network cannot represent an object rotated by $45^\circ$: any intermediate
orientation is approximated by the nearest discrete one, with an
\emph{aliasing} error that can reach $45^\circ$ for $C_4$. The remedy is to pass
to the continuous rotation group $\SO(2)$ and represent filters in a basis that
rotates analytically.

\subsection{Steerable filters}
\label{sec:groups:steerable-filters}

\begin{definition}[The group $\SO(2)$]
\label{def:groups:so2}
The special orthogonal group $\SO(2) = \{R_\theta : \theta \in [0, 2\pi)\}$
consists of all planar rotations about the origin, with
\begin{equation}
  R_\theta = \begin{pmatrix} \cos\theta & -\sin\theta \\
                             \sin\theta & \cos\theta \end{pmatrix}.
\end{equation}
\end{definition}

\begin{proposition}[Properties of $\SO(2)$]
\label{prop:groups:so2-properties}
The group $\SO(2)$ is \emph{compact} (topologically a circle $S^1$),
\emph{abelian} ($R_\alpha R_\beta = R_{\alpha+\beta} = R_\beta R_\alpha$), and
\emph{one-parameter} (parametrized by a single angle).
\end{proposition}

\begin{proof}
\emph{Compact:} the map $\theta \mapsto R_\theta$ is a homeomorphism of
$[0,2\pi)/(0 \sim 2\pi)$ onto $\SO(2)$, identifying it with $S^1$, which is
closed and bounded in $\R^{2 \times 2}$, hence compact by Heine--Borel.
\emph{Abelian:} $R_\alpha R_\beta$ is the rotation by $\alpha + \beta$, and
addition modulo $2\pi$ is commutative. \emph{One-parameter:} the conditions
$R^\top R = I$ and $\det R = 1$ impose three independent constraints on the four
entries of a $2 \times 2$ matrix, leaving $4 - 3 = 1$ degree of freedom.
\end{proof}

The circular structure makes harmonic analysis the natural tool. A
\emph{steerable filter} is one whose rotation can be synthesized from a finite
basis.

\begin{definition}[Steerable filter]
\label{def:groups:steerable-filter}
A filter $h: \R^2 \to \R$ is \emph{steerable} of order $N$ if there are basis
filters $\{h_0, \dots, h_{N-1}\}$ and interpolation functions $\{k_n(\theta)\}$
such that the rotated filter $h_\theta$ satisfies
\begin{equation}
  h_\theta(x,y) = \sum_{n=0}^{N-1} k_n(\theta)\, h_n(x,y).
\end{equation}
\end{definition}

Circular harmonics provide a natural steerable basis, because they transform by a
simple phase under rotation.

\begin{theorem}[Circular harmonic representation]
\label{thm:groups:circular-harmonics}
Every function $f: S^1 \to \C$ admits a Fourier series $f(\theta) = \sum_{k \in
\Z} c_k e^{ik\theta}$. Under a rotation, $g_\phi(\theta) = f(\theta + \phi)$ has
coefficients $c_k e^{ik\phi}$.
\end{theorem}

\begin{proof}
The expansion follows from completeness of the orthonormal system
$\{e^{ik\theta}\}_{k \in \Z}$ in $L^2(S^1)$, with $c_k = \frac{1}{2\pi}
\int_0^{2\pi} f(\theta) e^{-ik\theta}\,d\theta$. For the transformation,
\[
  \hat g_k = \frac{1}{2\pi}\int_0^{2\pi} f(\theta + \phi) e^{-ik\theta}\,d\theta
  = \frac{1}{2\pi}\int_0^{2\pi} f(\alpha) e^{-ik(\alpha - \phi)}\,d\alpha
  = e^{ik\phi} c_k,
\]
using the substitution $\alpha = \theta + \phi$ and periodicity.
\end{proof}

\begin{definition}[Circular harmonic filter]
\label{def:groups:circular-harmonic-filter}
A circular harmonic filter of frequency $k$ in polar coordinates $(r, \theta)$
is $\psi_k(r, \theta) = R_k(r)\, e^{ik\theta}$, with radial profile $R_k(r)$ and
angular part $e^{ik\theta}$.
\end{definition}

\Cref{fig:groups:steerable} illustrates the steerability of an oriented edge
filter: any rotation of the base filter is a fixed linear combination of two
basis filters.

\begin{figure}[h]
\centering
\begin{tikzpicture}[
    >=Stealth,
    every node/.style={font=\small},
    section label/.style={font=\normalsize\bfseries, text=gdlgray!70!black}
]
\newcommand{\steerablefilter}[4]{%
    \begin{scope}[shift={(#1,#2)}, rotate=#3]
        \fill[white, draw=gdlblue!60, thick] (0,0) -- (90:#4) arc (90:270:#4) -- cycle;
        \fill[gdlblue!30, draw=gdlblue!60, thick] (0,0) -- (270:#4) arc (270:450:#4) -- cycle;
        \draw[gdlblue!60, thick] (0,0) circle (#4);
        \draw[gdlgray!50, thin] (0,-#4) -- (0,#4);
    \end{scope}
}
\newcommand{\steerablefiltercolor}[5]{%
    \begin{scope}[shift={(#1,#2)}, rotate=#3]
        \fill[white, draw=#5, thick] (0,0) -- (90:#4) arc (90:270:#4) -- cycle;
        \fill[gdlblue!30, draw=#5, thick] (0,0) -- (270:#4) arc (270:450:#4) -- cycle;
        \draw[#5, thick] (0,0) circle (#4);
        \draw[gdlgray!50, thin] (0,-#4) -- (0,#4);
    \end{scope}
}
\node[section label, anchor=west] at (-1.2, 2.8) {Rotations of the base filter};
\steerablefilter{0}{1.2}{0}{0.7}
\node[below, text=gdlblue!70!black, font=\small] at (0, 0.35) {Base filter $\psi_0$};
\draw[->, thick, gdlgray!60] (1.0, 1.2) -- (1.8, 1.2);
\steerablefilter{2.8}{1.2}{0}{0.6}
\node[below, text=gdlpurple!70!black, font=\scriptsize] at (2.8, 0.45) {$0^\circ$};
\steerablefilter{4.6}{1.2}{45}{0.6}
\node[below, text=gdlpurple!70!black, font=\scriptsize] at (4.6, 0.45) {$45^\circ$};
\steerablefilter{6.4}{1.2}{90}{0.6}
\node[below, text=gdlpurple!70!black, font=\scriptsize] at (6.4, 0.45) {$90^\circ$};
\steerablefilter{8.2}{1.2}{135}{0.6}
\node[below, text=gdlpurple!70!black, font=\scriptsize] at (8.2, 0.45) {$135^\circ$};
\draw[gdlpurple!60, thick, decorate, decoration={brace, amplitude=5pt, mirror}]
    (2.1, 0.1) -- (8.9, 0.1);
\node[below=6pt, text=gdlpurple!70!black, font=\small] at (5.5, 0.05) {$\rho(r_\theta)\psi_0$};
\draw[->, gdlorange!70, thick, dashed] (3.5, 2.0) to[bend left=25] (4.5, 2.0);
\draw[->, gdlorange!70, thick, dashed] (5.3, 2.0) to[bend left=25] (6.3, 2.0);
\draw[->, gdlorange!70, thick, dashed] (7.1, 2.0) to[bend left=25] (8.1, 2.0);
\node[above, font=\tiny, text=gdlorange!70!black] at (4.0, 2.2) {$+45^\circ$};
\node[above, font=\tiny, text=gdlorange!70!black] at (5.8, 2.2) {$+45^\circ$};
\node[above, font=\tiny, text=gdlorange!70!black] at (7.6, 2.2) {$+45^\circ$};
\draw[gdlgray!30, thin] (-1.2, -0.7) -- (10.5, -0.7);
\node[section label, anchor=west] at (-1.2, -1.2) {Steerability};
\steerablefiltercolor{0.5}{-2.7}{35}{0.65}{gdlorange!70}
\node[below, text=gdlorange!70!black, font=\small] at (0.5, -3.6) {$\psi_\theta$};
\node[font=\Large, text=gdlgray!70!black] at (1.8, -2.7) {$=$};
\node[font=\normalsize, text=gdlred!70!black] at (2.7, -2.7) {$\cos\theta$};
\node[font=\large, text=gdlgray!70!black] at (3.4, -2.7) {$\times$};
\steerablefiltercolor{4.3}{-2.7}{0}{0.55}{gdlred!70}
\node[below, text=gdlred!70!black, font=\small] at (4.3, -3.45) {$\psi_0$};
\node[font=\Large, text=gdlgray!70!black] at (5.4, -2.7) {$+$};
\node[font=\normalsize, text=gdlgreen!70!black] at (6.3, -2.7) {$\sin\theta$};
\node[font=\large, text=gdlgray!70!black] at (7.0, -2.7) {$\times$};
\steerablefiltercolor{7.9}{-2.7}{90}{0.55}{gdlgreen!70}
\node[below, text=gdlgreen!70!black, font=\small] at (7.9, -3.45) {$\psi_{90}$};
\fill[gdlblue!8, rounded corners=4pt] (0.2, -4.4) rectangle (10.3, -5.4);
\draw[gdlblue!40, rounded corners=4pt, thick] (0.2, -4.4) rectangle (10.3, -5.4);
\node[font=\normalsize, text=gdlblue!70!black, anchor=west] at (0.6, -4.9)
    {Steerability:\quad $\displaystyle \rho(r_\theta)\psi = \sum_{k} a_k(\theta)\,\psi_k$};
\node[font=\scriptsize, text=gdlgray!70!black, anchor=west] at (7.0, -4.9)
    {$\{\psi_k\}$: filter basis};
\end{tikzpicture}
\caption[Steerable filters]{Steerable filters: any rotation of the base filter
$\psi_0$ is a fixed linear combination of a finite basis of filters. This
property delivers equivariance to continuous rotations with finitely many
parameters.}
\label{fig:groups:steerable}
\end{figure}

In practice a steerable convolution stores learnable coefficients that combine a
fixed circular-harmonic basis (a Gaussian radial envelope times $\cos k\theta$
and $\sin k\theta$ for $k = 0, 1, \dots, F$) into the convolution kernel.
Equivariant nonlinearities preserve the irreducible structure: the trivial
($k=0$) channels use ordinary pointwise activations, while nontrivial channels
use a norm-nonlinearity that gates the magnitude of the feature while preserving
its phase. Steerable networks following the framework of
\citet{weiler2019general} and the early steerable filter learning of
\citet{weiler2018learningsteerable} outperform finite discretizations on rotated
benchmarks while using fewer parameters, because the harmonic basis represents
continuous rotational symmetry exactly rather than by discrete copies.

\subsection{Irreducible representations and geometric fields}
\label{sec:groups:irreps-fields}

Steerable networks operate not on scalar channels but on \emph{geometric
fields} that transform under irreducible representations of $\SO(2)$.

\begin{definition}[Geometric field]
\label{def:groups:geometric-field}
A geometric field of type $\rho$ is a function $f: \R^2 \to V$ where $V$ carries
a representation $\rho: \SO(2) \to \GL(V)$ such that
\begin{equation}
  (g \cdot f)(x) = \rho(g)\, f(g^{-1} x) \qquad \forall g \in \SO(2).
\end{equation}
\end{definition}

\begin{definition}[Irreducible representations of $\SO(2)$]
\label{def:groups:so2-irreps}
The irreducible representations of $\SO(2)$ are indexed by integers $k \in \Z$:
$\rho_k(R_\theta) = e^{ik\theta}$ for $k \neq 0$, and the trivial representation
$\rho_0(R_\theta) = 1$. Over the reals, frequencies $\pm k$ combine into the
two-dimensional representation
\begin{equation}
  \rho_k^{\text{real}}(R_\theta) = \begin{pmatrix}
    \cos(k\theta) & -\sin(k\theta) \\
    \sin(k\theta) & \cos(k\theta) \end{pmatrix}.
\end{equation}
\end{definition}

Steerable networks excel where exact orientation is informative: fiber analysis
in medical imaging, fingerprint ridge detection, and crack and defect detection
in materials. They form the natural bridge to the three-dimensional rotation
group $\SO(3)$, to which we turn next.

\section{Spherical CNNs and $\SO(3)$ equivariance}
\label{sec:groups:spherical}

Many real-world signals live on the sphere or carry full three-dimensional
rotational symmetry: $360^\circ$ omnidirectional images are defined on $S^2$;
global climate fields live on the (topologically spherical) surface of the
Earth; the cosmic microwave background is a signal on the celestial sphere.
Projecting such data to a rectangular grid introduces area distortion near the
poles, artificial discontinuities at the seams, and a loss of equivariance,
since spherical rotations become complicated transformations in the projected
domain. The principled alternative is to build convolutions equivariant to the
rotation group $\SO(3)$.

\begin{definition}[The group $\SO(3)$]
\label{def:groups:so3}
$\SO(3) = \{R \in \R^{3 \times 3} : R^\top R = I,\ \det R = 1\}$ is the group of
orientation-preserving rotations of $\R^3$. Each element may be parametrized by
Euler angles $(\alpha, \beta, \gamma)$ or by a rotation axis $\mathbf{n} \in S^2$
and angle $\theta \in [0, 2\pi)$.
\end{definition}

\begin{proposition}[Properties of $\SO(3)$]
\label{prop:groups:so3-properties}
The group $\SO(3)$ is \emph{non-abelian}, \emph{compact} (topologically the real
projective space $\R P^3$), and \emph{three-parameter}.
\end{proposition}

\begin{proof}
\emph{Non-abelian:} let $R_x, R_z$ be $90^\circ$ rotations about the $x$- and
$z$-axes; then $R_x R_z (1,0,0)^\top = (0,0,1)^\top$ but $R_z R_x (1,0,0)^\top =
(0,1,0)^\top$, so $R_x R_z \neq R_z R_x$. \emph{Compact:} $\SO(3)$ is closed in
$\R^9$ and bounded ($\norm{R}_F^2 = \operatorname{tr}(R^\top R) = 3$), hence
compact; the homeomorphism with $\R P^3$ comes from the axis--angle
parametrization, identifying $(\pi, \mathbf{n})$ with $(\pi, -\mathbf{n})$.
\emph{Three-parameter:} the condition $R^\top R = I$ imposes six independent
constraints on the nine entries, so $\dim \SO(3) = 9 - 6 = 3$.
\end{proof}

\subsection{Spherical harmonics}
\label{sec:groups:spherical-harmonics}

Harmonic analysis on $S^2$ provides the basis for $\SO(3)$-equivariant
convolution.

\begin{definition}[Spherical harmonics]
\label{def:groups:spherical-harmonics}
The spherical harmonics $Y_\ell^m(\theta, \phi)$ form a complete orthonormal
basis of $L^2(S^2)$:
\begin{equation}
  Y_\ell^m(\theta, \phi) = \sqrt{\frac{2\ell+1}{4\pi}\,
    \frac{(\ell - |m|)!}{(\ell + |m|)!}}\,
    P_\ell^{|m|}(\cos\theta)\, e^{im\phi},
\end{equation}
where $P_\ell^m$ are the associated Legendre polynomials, $\ell \geq 0$ is the
degree, and $|m| \leq \ell$ the order.
\end{definition}

Every $f \in L^2(S^2)$ expands as $f(\theta, \phi) = \sum_{\ell=0}^{\infty}
\sum_{m=-\ell}^{\ell} f_\ell^m\, Y_\ell^m(\theta, \phi)$ with coefficients
$f_\ell^m = \int_{S^2} f(\mathbf{x})\, \overline{Y_\ell^m(\mathbf{x})}\,
d\mathbf{x}$. \Cref{fig:groups:spherical-harmonics} displays the low-degree
harmonics. Their defining property is their behavior under rotation.

\begin{theorem}[Transformation of spherical harmonics]
\label{thm:groups:sh-transformation}
Under a rotation $R \in \SO(3)$,
\begin{equation}
  Y_\ell^m(R^{-1}\mathbf{x}) = \sum_{m'=-\ell}^{\ell} D^{(\ell)}_{m'm}(R)\,
    Y_\ell^{m'}(\mathbf{x}),
\end{equation}
where $D^{(\ell)}(R)$ is the $(2\ell+1)\times(2\ell+1)$ irreducible Wigner matrix
of $\SO(3)$.
\end{theorem}

\begin{proof}
The harmonics $\{Y_\ell^m\}_{m=-\ell}^{\ell}$ span the space $\mathcal{H}_\ell$
of homogeneous harmonic polynomials of degree $\ell$ restricted to $S^2$. This
space is $\SO(3)$-invariant, because composition with a rotation preserves both
homogeneity and harmonicity (the Laplacian commutes with rotations). Since
$\mathcal{H}_\ell$ has dimension $2\ell+1$ and is invariant, $Y_\ell^m \circ
R^{-1}$ is a linear combination of the basis, with coefficients $D^{(\ell)}_{m'm}
(R) = \int_{S^2} Y_\ell^m(R^{-1}\mathbf{x})\, \overline{Y_\ell^{m'}(\mathbf{x})}
\, d\mathbf{x}$. Irreducibility of $D^{(\ell)}$ follows because
$\mathcal{H}_\ell$ has no proper $\SO(3)$-invariant subspace.
\end{proof}

\begin{figure}[h]
\centering
\begin{tikzpicture}[every node/.style={font=\small}]
  \def\sphR{1}
  \def\colsep{2.8}
  \def\rowsep{3.2}
  \node[font=\large\bfseries, text=gdlblue!70!black]
    at ({2*\colsep}, 2.4) {Spherical harmonics $Y_\ell^m(\theta,\varphi)$};
  \foreach \m/\col in {-2/0, -1/1, 0/2, 1/3, 2/4} {
    \node[font=\footnotesize, text=gdlpurple!70!black] at ({\col*\colsep}, 1.5) {$m = \m$};
  }
  \node[font=\footnotesize, text=gdlorange!70!black] at (-1.6, 0)          {$\ell = 0$};
  \node[font=\footnotesize, text=gdlorange!70!black] at (-1.6, -\rowsep)   {$\ell = 1$};
  \node[font=\footnotesize, text=gdlorange!70!black] at (-1.6, -2*\rowsep) {$\ell = 2$};
  % l=0,m=0
  \begin{scope}[shift={(2*\colsep, 0)}]
    \fill[gdlblue!20] (0,0) circle (\sphR);
    \draw[gdlgray!60, thick] (0,0) circle (\sphR);
    \draw[gdlgray!40, dashed] ({\sphR},0) arc[start angle=0, end angle=180, x radius=\sphR, y radius=0.3*\sphR];
    \draw[gdlgray!50] ({-\sphR},0) arc[start angle=180, end angle=360, x radius=\sphR, y radius=0.3*\sphR];
    \node[font=\scriptsize, text=gdlblue!70!black] at (0, -\sphR - 0.4) {$Y_0^0$};
  \end{scope}
  % l=1,m=-1
  \begin{scope}[shift={(1*\colsep, -\rowsep)}]
    \fill[gdlblue!30] (90:\sphR) arc (90:270:\sphR) -- cycle;
    \fill[gdlred!30]  (90:\sphR) arc (90:-90:\sphR) -- cycle;
    \draw[gdlgray!60, thick] (0,0) circle (\sphR);
    \draw[gdlgray!40, dashed] (0,\sphR) -- (0,-\sphR);
    \draw[gdlgray!40, dashed] ({\sphR},0) arc[start angle=0, end angle=180, x radius=\sphR, y radius=0.3*\sphR];
    \draw[gdlgray!50] ({-\sphR},0) arc[start angle=180, end angle=360, x radius=\sphR, y radius=0.3*\sphR];
    \node[font=\scriptsize, text=gdlblue!70!black] at (-0.5, 0.5) {$+$};
    \node[font=\scriptsize, text=gdlred!70!black]  at ( 0.5, 0.5) {$-$};
    \node[font=\scriptsize, text=gdlpurple!70!black] at (0, -\sphR - 0.4) {$Y_1^{-1}$};
  \end{scope}
  % l=1,m=0
  \begin{scope}[shift={(2*\colsep, -\rowsep)}]
    \fill[gdlblue!30] (-\sphR,0) arc (180:0:\sphR) -- cycle;
    \fill[gdlred!30]  (-\sphR,0) arc (180:360:\sphR) -- cycle;
    \draw[gdlgray!60, thick] (0,0) circle (\sphR);
    \draw[gdlgray!50] (-\sphR,0) -- (\sphR,0);
    \draw[gdlgray!40, dashed] ({\sphR},0) arc[start angle=0, end angle=180, x radius=\sphR, y radius=0.3*\sphR];
    \draw[gdlgray!50] ({-\sphR},0) arc[start angle=180, end angle=360, x radius=\sphR, y radius=0.3*\sphR];
    \node[font=\scriptsize, text=gdlblue!70!black] at (0,  0.5) {$+$};
    \node[font=\scriptsize, text=gdlred!70!black]  at (0, -0.5) {$-$};
    \node[font=\scriptsize, text=gdlpurple!70!black] at (0, -\sphR - 0.4) {$Y_1^{0}$};
  \end{scope}
  % l=1,m=1
  \begin{scope}[shift={(3*\colsep, -\rowsep)}]
    \fill[gdlred!30]  (90:\sphR) arc (90:270:\sphR) -- cycle;
    \fill[gdlblue!30] (90:\sphR) arc (90:-90:\sphR) -- cycle;
    \draw[gdlgray!60, thick] (0,0) circle (\sphR);
    \draw[gdlgray!40, dashed] (0,\sphR) -- (0,-\sphR);
    \draw[gdlgray!40, dashed] ({\sphR},0) arc[start angle=0, end angle=180, x radius=\sphR, y radius=0.3*\sphR];
    \draw[gdlgray!50] ({-\sphR},0) arc[start angle=180, end angle=360, x radius=\sphR, y radius=0.3*\sphR];
    \node[font=\scriptsize, text=gdlred!70!black]  at (-0.5, 0.5) {$-$};
    \node[font=\scriptsize, text=gdlblue!70!black] at ( 0.5, 0.5) {$+$};
    \node[font=\scriptsize, text=gdlpurple!70!black] at (0, -\sphR - 0.4) {$Y_1^{1}$};
  \end{scope}
  % l=2,m=-2
  \begin{scope}[shift={(0*\colsep, -2*\rowsep)}]
    \clip (0,0) circle (\sphR);
    \fill[gdlblue!30] (0,0) -- (90:\sphR)  arc (90:0:\sphR)    -- cycle;
    \fill[gdlred!30]  (0,0) -- (0:\sphR)   arc (0:-90:\sphR)   -- cycle;
    \fill[gdlblue!30] (0,0) -- (-90:\sphR) arc (-90:-180:\sphR) -- cycle;
    \fill[gdlred!30]  (0,0) -- (180:\sphR) arc (180:90:\sphR)  -- cycle;
    \draw[gdlgray!40, dashed] (-\sphR,0) -- (\sphR,0);
    \draw[gdlgray!40, dashed] (0,-\sphR) -- (0,\sphR);
    \draw[gdlgray!60, thick] (0,0) circle (\sphR);
    \node[font=\scriptsize, text=gdlpurple!70!black] at (0, -\sphR - 0.4) {$Y_2^{-2}$};
  \end{scope}
  % l=2,m=-1
  \begin{scope}[shift={(1*\colsep, -2*\rowsep)}]
    \clip (0,0) circle (\sphR);
    \fill[gdlblue!30] (0,0) -- (45:\sphR)  arc (45:135:\sphR)  -- cycle;
    \fill[gdlred!30]  (0,0) -- (135:\sphR) arc (135:225:\sphR) -- cycle;
    \fill[gdlblue!30] (0,0) -- (225:\sphR) arc (225:315:\sphR) -- cycle;
    \fill[gdlred!30]  (0,0) -- (315:\sphR) arc (315:405:\sphR) -- cycle;
    \draw[gdlgray!60, thick] (0,0) circle (\sphR);
    \node[font=\scriptsize, text=gdlpurple!70!black] at (0, -\sphR - 0.4) {$Y_2^{-1}$};
  \end{scope}
  % l=2,m=0
  \begin{scope}[shift={(2*\colsep, -2*\rowsep)}]
    \clip (0,0) circle (\sphR);
    \fill[gdlblue!30] (-\sphR, {0.33*\sphR}) rectangle (\sphR, \sphR);
    \fill[gdlred!30]  (-\sphR, {-0.33*\sphR}) rectangle (\sphR, {0.33*\sphR});
    \fill[gdlblue!30] (-\sphR, -\sphR) rectangle (\sphR, {-0.33*\sphR});
    \draw[gdlgray!60, thick] (0,0) circle (\sphR);
    \node[font=\scriptsize, text=gdlpurple!70!black] at (0, -\sphR - 0.4) {$Y_2^{0}$};
  \end{scope}
  % l=2,m=1
  \begin{scope}[shift={(3*\colsep, -2*\rowsep)}]
    \clip (0,0) circle (\sphR);
    \fill[gdlred!30]  (0,0) -- (45:\sphR)  arc (45:135:\sphR)  -- cycle;
    \fill[gdlblue!30] (0,0) -- (135:\sphR) arc (135:225:\sphR) -- cycle;
    \fill[gdlred!30]  (0,0) -- (225:\sphR) arc (225:315:\sphR) -- cycle;
    \fill[gdlblue!30] (0,0) -- (315:\sphR) arc (315:405:\sphR) -- cycle;
    \draw[gdlgray!60, thick] (0,0) circle (\sphR);
    \node[font=\scriptsize, text=gdlpurple!70!black] at (0, -\sphR - 0.4) {$Y_2^{1}$};
  \end{scope}
  % l=2,m=2
  \begin{scope}[shift={(4*\colsep, -2*\rowsep)}]
    \clip (0,0) circle (\sphR);
    \fill[gdlred!30]  (0,0) -- (90:\sphR)  arc (90:0:\sphR)    -- cycle;
    \fill[gdlblue!30] (0,0) -- (0:\sphR)   arc (0:-90:\sphR)   -- cycle;
    \fill[gdlred!30]  (0,0) -- (-90:\sphR) arc (-90:-180:\sphR) -- cycle;
    \fill[gdlblue!30] (0,0) -- (180:\sphR) arc (180:90:\sphR)  -- cycle;
    \draw[gdlgray!60, thick] (0,0) circle (\sphR);
    \node[font=\scriptsize, text=gdlpurple!70!black] at (0, -\sphR - 0.4) {$Y_2^{2}$};
  \end{scope}
  \begin{scope}[shift={(0, {-2*\rowsep - 2.3})}]
    \fill[gdlblue!30] ({0.8*\colsep}, 0) rectangle ++(0.4, 0.4);
    \draw[gdlgray!60] ({0.8*\colsep}, 0) rectangle ++(0.4, 0.4);
    \node[font=\scriptsize, anchor=west, text=gdlblue!70!black]
      at ({0.8*\colsep + 0.55}, 0.2) {Positive region ($Y_\ell^m > 0$)};
    \fill[gdlred!30] ({2.5*\colsep}, 0) rectangle ++(0.4, 0.4);
    \draw[gdlgray!60] ({2.5*\colsep}, 0) rectangle ++(0.4, 0.4);
    \node[font=\scriptsize, anchor=west, text=gdlred!70!black]
      at ({2.5*\colsep + 0.55}, 0.2) {Negative region ($Y_\ell^m < 0$)};
  \end{scope}
  \node[font=\scriptsize, text=gdlgray!80!black, text width=9cm, align=center]
    at ({2*\colsep}, {-2*\rowsep - 3.3})
    {$\ell$ = degree (frequency),\quad $m$ = order (orientation)};
\end{tikzpicture}
\caption[Spherical harmonics]{Spherical harmonics $Y_\ell^m(\theta, \varphi)$
for degrees $\ell = 0, 1, 2$. The degree $\ell$ controls the angular frequency
(number of sign changes) and the order $m$ the orientation of the nodal pattern.
These functions form the basis of the $(2\ell+1)$-dimensional irreducible
representations of $\SO(3)$.}
\label{fig:groups:spherical-harmonics}
\end{figure}

\subsection{Spherical convolution}
\label{sec:groups:spherical-convolution}

Spherical convolutional networks, introduced by \citet{cohen2018sphericalcnns},
process functions on $S^2 = \SO(3)/\SO(2)$ directly through the group structure
of $\SO(3)$.

\begin{definition}[$\SO(3)$ convolution]
\label{def:groups:so3-convolution}
For $f: \SO(3) \to \R^{C_{\text{in}}}$ and kernel $\psi: \SO(3) \to
\R^{C_{\text{out}} \times C_{\text{in}}}$,
\begin{equation}
  (f \star \psi)(g) = \int_{\SO(3)} f(h)\, \psi(g^{-1} h)\, dh,
\end{equation}
where $dh$ is the Haar measure on $\SO(3)$.
\end{definition}

\begin{theorem}[Equivariance of spherical convolution]
\label{thm:groups:spherical-conv-equivariance}
The $\SO(3)$ convolution is equivariant: if $(g \cdot f)(x) = f(g^{-1} x)$ then
$g \cdot (f \star \psi) = (g \cdot f) \star \psi$.
\end{theorem}

\begin{proof}
Compute, for $g \in \SO(3)$,
\begin{align}
  [(g \cdot f) \star \psi](x)
  &= \int_{\SO(3)} (g \cdot f)(h)\, \psi(x^{-1} h)\, dh
   = \int_{\SO(3)} f(g^{-1} h)\, \psi(x^{-1} h)\, dh.
\end{align}
Substituting $h' = g^{-1} h$ (so $h = g h'$) and using left-invariance of the
Haar measure ($dh = dh'$),
\begin{align}
  &= \int_{\SO(3)} f(h')\, \psi(x^{-1} g h')\, dh'
   = \int_{\SO(3)} f(h')\, \psi\big((g^{-1} x)^{-1} h'\big)\, dh'
   = (f \star \psi)(g^{-1} x),
\end{align}
which is exactly $[g \cdot (f \star \psi)](x)$.
\end{proof}

Efficient implementations use a spherical fast Fourier transform: the convolution
theorem extends to the sphere, reducing convolution to matrix multiplication
within each degree-$\ell$ block, $\widehat{(f \star \psi)}_\ell^m = \sum_{m'}
\hat f_\ell^{m'} \hat\psi_\ell^{m'm}$. The complexity is governed by the maximum
degree $L$: $O(L^2)$ coefficients per spherical location and $O(L^3)$ for the
transform itself. The construction extends to volumetric data by combining
spherical harmonic analysis with radial profiles, as in the 3D steerable
networks of \citet{weiler2018steerable}, and to unstructured spherical grids
\citep{cohen2021spherical}. The progression from standard convolutional networks
to discrete group convolutional networks, to $\SO(2)$ steerable networks, to
$\SO(3)$ spherical networks traces a steady increase in geometric awareness,
each level matching the architecture more closely to the intrinsic symmetry of
the data.

\section{Non-compact groups: $\Eucl(3)$ and $\SE(3)$ equivariance}
\label{sec:groups:noncompact}

Point clouds, molecules, and particle systems carry symmetries that include
\emph{translations} as well as rotations. The relevant groups are the Euclidean
group $\Eucl(n)$ and the special Euclidean group $\SE(n)$, which are
\emph{non-compact}: the translation subgroup $\R^n$ is unbounded, so the Haar
measure is infinite and direct integration over the group is no longer feasible.
Equivariant architectures must therefore be built differently.

\begin{definition}[Euclidean and special Euclidean groups]
\label{def:groups:euclidean}
The \emph{Euclidean group} $\Eucl(n) = \R^n \rtimes \OO(n)$ is the group of all
isometries of $\R^n$. Each element is a pair $(R, \mathbf{t})$ with $R \in
\OO(n)$ and $\mathbf{t} \in \R^n$, acting by $g \cdot \mathbf{x} = R\mathbf{x} +
\mathbf{t}$, with composition $(R_1, \mathbf{t}_1)\circ(R_2, \mathbf{t}_2) =
(R_1 R_2,\, R_1 \mathbf{t}_2 + \mathbf{t}_1)$. The \emph{special Euclidean group}
$\SE(n) = \R^n \rtimes \SO(n)$ is the orientation-preserving subgroup.
\end{definition}

\begin{proposition}[Properties of $\Eucl(n)$]
\label{prop:groups:en-properties}
$\Eucl(n)$ is \emph{non-compact} (its translation subgroup $\R^n$ is unbounded),
a \emph{Lie group} of dimension $n + \tfrac{n(n-1)}{2}$, and a \emph{semidirect
product} in which $\R^n$ is a normal subgroup.
\end{proposition}

\begin{proof}
\emph{Non-compact:} pure translations $\{(I, \mathbf{t})\} \cong \R^n$ are
unbounded. \emph{Lie group:} $\Eucl(n)$ embeds as a closed subgroup of
$\GL(n+1, \R)$ via $(R, \mathbf{t}) \mapsto \left(\begin{smallmatrix} R &
\mathbf{t} \\ 0 & 1 \end{smallmatrix}\right)$, of dimension $n + \dim \OO(n) = n
+ \tfrac{n(n-1)}{2}$. \emph{Semidirect product:} conjugating a pure translation,
$(R, \mathbf{s})(I, \mathbf{t})(R, \mathbf{s})^{-1} = (I, R\mathbf{t})$, remains
a translation, so $\R^n$ is normal and the composition law is exactly that of
$\R^n \rtimes \OO(n)$.
\end{proof}

The crucial structural fact for architecture design is that a finite-dimensional
representation in which translations act trivially factors through the rotation
part.

\begin{proposition}[Representation decomposition]
\label{prop:groups:rep-decomposition}
In any finite-dimensional representation of $\Eucl(n) = \R^n \rtimes \OO(n)$ in
which the translations act trivially, the representation factors through the
quotient $\Eucl(n)/\R^n \cong \OO(n)$: $\rho(R, \mathbf{t}) = \rho_{\OO(n)}(R)$.
The analogous statement holds for $\SE(n)$ through $\SO(n)$
\citep{DBLP:journals/corr/abs-2104-13478}.
\end{proposition}

\begin{proof}
Let $\rho: \Eucl(n) \to \GL(V)$ satisfy $\rho(I, \mathbf{t}) = \operatorname{Id}$
for all $\mathbf{t}$. Since $(R, \mathbf{t}) = (R, \mathbf{0})\cdot(I, R^{-1}
\mathbf{t})$, we have $\rho(R, \mathbf{t}) = \rho(R, \mathbf{0})$. Setting
$\rho_{\OO(n)}(R) := \rho(R, \mathbf{0})$ gives a homomorphism, because
$\rho_{\OO(n)}(R_1 R_2) = \rho(R_1 R_2, \mathbf{0}) = \rho(R_1, \mathbf{0})
\rho(R_2, \mathbf{0})$. Thus $\rho$ factors through the canonical projection
$\pi: \Eucl(n) \to \Eucl(n)/\R^n \cong \OO(n)$.
\end{proof}

This factorization motivates architectures that separate \emph{invariant}
information (scalar features) from \emph{equivariant} information (vector
coordinates), and yields two families of models: graph networks built on
invariant scalars, and tensor-field networks built on full $\SO(3)$
representations.

\subsection{Invariant messages: \texorpdfstring{$\Eucl(n)$}{E(n)}-equivariant graph networks}
\label{sec:groups:egnn}

The $\Eucl(n)$-equivariant graph neural networks of
\citet{satorras2021enequivariant} (EGNNs) build messages from $\Eucl(n)$-invariant
quantities while updating coordinates through vector differences that preserve
equivariance. Each node carries scalar features $h_i \in \R^d$ (invariant under
$\Eucl(n)$) and coordinates $\mathbf{x}_i \in \R^n$ (equivariant under
$\Eucl(n)$), updated in three stages.

\paragraph{Stage 1: invariant messages.}
Messages depend only on $\Eucl(n)$-invariant quantities:
\begin{equation}
  \mathbf{m}_{ij} = \phi_e\big(h_i, h_j, \norm{\mathbf{x}_i - \mathbf{x}_j}^2,
    a_{ij}\big),
  \label{eq:groups:egnn-message}
\end{equation}
where $h_i, h_j$ are scalar features, $\norm{\mathbf{x}_i - \mathbf{x}_j}^2$ is
the squared Euclidean distance, $a_{ij}$ are optional edge attributes, and
$\phi_e$ is a multilayer perceptron. The squared distance is $\Eucl(n)$-invariant:
for any $(R, \mathbf{t})$,
\begin{equation}
  \norm{(R\mathbf{x}_i + \mathbf{t}) - (R\mathbf{x}_j + \mathbf{t})}^2
  = \norm{R(\mathbf{x}_i - \mathbf{x}_j)}^2
  = \norm{\mathbf{x}_i - \mathbf{x}_j}^2,
\end{equation}
since $R$ is orthogonal and translations cancel.

\paragraph{Stage 2: equivariant coordinate update.}
Coordinates are updated by weighted vector differences:
\begin{equation}
  \mathbf{x}_i' = \mathbf{x}_i + C \sum_{j \in \mathcal{N}(i)}
    (\mathbf{x}_i - \mathbf{x}_j)\, \phi_x(\mathbf{m}_{ij}),
  \label{eq:groups:egnn-coord}
\end{equation}
where $\phi_x$ is a scalar-valued multilayer perceptron, $C$ a normalization
constant, and $\mathcal{N}(i)$ the neighborhood of node $i$.

\begin{proposition}[Equivariance of the coordinate update]
\label{prop:groups:egnn-equivariance}
The update \Cref{eq:groups:egnn-coord} is $\Eucl(n)$-equivariant: applying
$(R, \mathbf{t})$ to all coordinates yields $\mathbf{x}_i' \mapsto R\mathbf{x}_i'
+ \mathbf{t}$.
\end{proposition}

\begin{proof}
With transformed coordinates $\tilde{\mathbf{x}}_i = R\mathbf{x}_i + \mathbf{t}$,
the messages $\mathbf{m}_{ij}$ are unchanged (built from invariants), so
\begin{align}
  \tilde{\mathbf{x}}_i'
  &= R\mathbf{x}_i + \mathbf{t} + C \sum_{j} \big((R\mathbf{x}_i + \mathbf{t})
     - (R\mathbf{x}_j + \mathbf{t})\big)\, \phi_x(\mathbf{m}_{ij}) \\
  &= R\mathbf{x}_i + \mathbf{t} + C \sum_{j} R(\mathbf{x}_i - \mathbf{x}_j)\,
     \phi_x(\mathbf{m}_{ij})
   = R\Big(\mathbf{x}_i + C \sum_{j} (\mathbf{x}_i - \mathbf{x}_j)\,
     \phi_x(\mathbf{m}_{ij})\Big) + \mathbf{t}
   = R\mathbf{x}_i' + \mathbf{t}. \qedhere
\end{align}
\end{proof}

\paragraph{Stage 3: invariant feature update.}
Scalar features are updated from aggregated messages, $\mathbf{m}_i = \sum_{j \in
\mathcal{N}(i)} \mathbf{m}_{ij}$ and $h_i' = \phi_h(h_i, \mathbf{m}_i)$, which is
trivially invariant because it operates only on invariant quantities. The
architecture thus guarantees $h_i'(g \cdot \mathbf{X}) = h_i'(\mathbf{X})$
(invariance) and $\mathbf{x}_i'(g \cdot \mathbf{X}) = g \cdot
\mathbf{x}_i'(\mathbf{X})$ (equivariance) for every $g = (R, \mathbf{t})$.

\begin{example}[EGNN message passing for a small molecule]
\label{ex:groups:egnn-molecular}
Consider three atoms at $\mathbf{x}_1 = (0,0,0)$, $\mathbf{x}_2 = (1.5,0,0)$,
$\mathbf{x}_3 = (0.75, 1.3, 0)$ with scalar features $h_1 = [6.0]$ (carbon),
$h_2 = [1.0]$ (hydrogen), $h_3 = [8.0]$ (oxygen). The squared distances are
$d_{12}^2 = 2.25$, $d_{13}^2 = 0.75^2 + 1.3^2 = 2.2525$, $d_{23}^2 = 2.2525$.
Messages are $\mathbf{m}_{12} = \phi_e(h_1, h_2, d_{12}^2) = \MLP([6.0, 1.0,
2.25])$, and the coordinate update is $\mathbf{x}_1' = \mathbf{x}_1 +
\tfrac{1}{2}[(\mathbf{x}_1 - \mathbf{x}_2)\phi_x(\mathbf{m}_{12}) + (\mathbf{x}_1
- \mathbf{x}_3)\phi_x(\mathbf{m}_{13})]$. Rotating the whole molecule rotates the
updated coordinates identically, by \Cref{prop:groups:egnn-equivariance}.
\end{example}

EGNNs are highly scalable and suited to large systems (thousands of atoms), at
the cost of expressivity: they use only scalar features and position vectors, so
they cannot represent higher-order directional interactions. Incorporating
additional invariants --- triplet angles $\theta_{ijk}$, dihedral angles
$\phi_{ijkl}$ (which are $\SE(3)$-invariant but change sign under reflection, and
so distinguish chirality), oriented volumes, or physical quantities such as
Coulomb and Lennard-Jones potentials --- improves performance without breaking
the symmetry \citep{tholke2022geometric}. For full $\Eucl(3)$ equivariance,
including reflections, one uses $|\phi_{ijkl}|$ instead of the signed dihedral
angle.

\subsection{Tensor features: Tensor Field Networks and $\SE(3)$-Transformers}
\label{sec:groups:tfn}

A more expressive family represents node features as full tensors transforming
under the irreducible representations of $\SO(3)$.

\begin{definition}[Irreducible representations of $\SO(3)$]
\label{def:groups:so3-irreps}
The irreducible representations of $\SO(3)$ are indexed by non-negative integers
$\ell = 0, 1, 2, \dots$ (the degree, or angular momentum). The degree-$\ell$
representation has dimension $2\ell + 1$, with matrix elements, in Euler angles
$(\alpha, \beta, \gamma)$,
\begin{equation}
  D^{(\ell)}_{mm'}(\alpha, \beta, \gamma)
  = e^{-im\alpha}\, d^{(\ell)}_{mm'}(\beta)\, e^{-im'\gamma},
\end{equation}
where $d^{(\ell)}_{mm'}$ are the small Wigner matrices.
\end{definition}

Low degrees have direct geometric meaning: $\ell = 0$ are scalars (invariant),
$\ell = 1$ are vectors ($\mathbf{v}' = R\mathbf{v}$), and $\ell = 2$ are
symmetric traceless second-order tensors. A degree-$\ell$ \emph{feature tensor}
$\mathbf{f}_i^{(\ell)} \in \R^{2\ell+1}$ at node $i$ transforms as
$(\mathbf{f}_i^{(\ell)})' = D^{(\ell)}(R)\, \mathbf{f}_i^{(\ell)}$, and each node
holds several copies (multiplicities) of each type, $\mathbf{F}_i =
\bigoplus_{\ell=0}^{L_{\max}} \bigoplus_{k=1}^{n_\ell}
\mathbf{f}_{i,k}^{(\ell)}$.

\Cref{fig:groups:so3} recalls the elementary rotation matrices and the basic
properties of $\SO(3)$. Combining features of different degrees requires the
Clebsch--Gordan tensor product.

\begin{figure}[h]
\centering
\begin{tikzpicture}[>=Stealth, scale=1.1]
\draw[thick, ->] (0,0,0) -- (3,0,0) node[right] {$x$};
\draw[thick, ->] (0,0,0) -- (0,3,0) node[above] {$y$};
\draw[thick, ->] (0,0,0) -- (0,0,3) node[below left] {$z$};
\draw[gdlgray!40, dashed] (0,0) circle (2cm);
\draw[->, very thick, gdlblue!70] (0,0,0) -- (2,0,0);
\fill[gdlblue!70] (2,0,0) circle (3pt);
\node[gdlblue!70, right] at (2.1,0,0) {$\mathbf{v}$};
\draw[->, very thick, gdlred!70] (0,0,0) -- ({2*cos(45)},{2*sin(45)},0);
\fill[gdlred!70] ({2*cos(45)},{2*sin(45)},0) circle (3pt);
\node[gdlred!70, above right] at ({2*cos(45)},{2*sin(45)},0) {$R_z(\theta)\mathbf{v}$};
\draw[thick, gdlorange!70, ->] (1.2,0,0) arc (0:45:1.2cm);
\node[gdlorange!70] at (1, 0.5, 0) {$\theta$};
\node[font=\small, text width=7cm] at (7, 0.8) {
    \textbf{Elementary rotations:}\\[0.4em]
    $R_x(\theta) = \begin{psmallmatrix} 1 & 0 & 0 \\ 0 & \cos\theta & -\sin\theta \\ 0 & \sin\theta & \cos\theta \end{psmallmatrix}$,\quad
    $R_y(\theta) = \begin{psmallmatrix} \cos\theta & 0 & \sin\theta \\ 0 & 1 & 0 \\ -\sin\theta & 0 & \cos\theta \end{psmallmatrix}$,\quad
    $R_z(\theta) = \begin{psmallmatrix} \cos\theta & -\sin\theta & 0 \\ \sin\theta & \cos\theta & 0 \\ 0 & 0 & 1 \end{psmallmatrix}$
};
\node[font=\small, text width=6cm] at (7, -2.4) {
    \textbf{Properties of $\SO(3)$:}
    \begin{itemize}
        \item Compact Lie group, dimension $3$
        \item $\det R = 1$, $R^\top R = I$
        \item Non-abelian: $R_x R_y \neq R_y R_x$
        \item Topology: $\R P^3$
    \end{itemize}
};
\end{tikzpicture}
\caption[Rotations in $\SO(3)$]{Rotations in three dimensions. Any $R \in \SO(3)$
is a rotation about some axis $\mathbf{n}$ by an angle $\theta$, and decomposes
into elementary rotations about the coordinate axes.}
\label{fig:groups:so3}
\end{figure}

\begin{theorem}[Clebsch--Gordan decomposition]
\label{thm:groups:clebsch-gordan}
The tensor product of two irreducible representations of $\SO(3)$ decomposes as a
direct sum of irreducibles:
\begin{equation}
  \rho^{(\ell_1)} \otimes \rho^{(\ell_2)}
  = \bigoplus_{\ell = |\ell_1 - \ell_2|}^{\ell_1 + \ell_2} \rho^{(\ell)}.
\end{equation}
The Clebsch--Gordan coefficients $C^{\ell m}_{\ell_1 m_1, \ell_2 m_2}$ implement
the decomposition:
\begin{equation}
  (f^{(\ell_1)} \otimes g^{(\ell_2)})^{(\ell)}_m
  = \sum_{m_1, m_2} C^{\ell m}_{\ell_1 m_1, \ell_2 m_2}\,
    f^{(\ell_1)}_{m_1}\, g^{(\ell_2)}_{m_2}.
\end{equation}
\end{theorem}

The selection rule $|\ell_1 - \ell_2| \leq \ell \leq \ell_1 + \ell_2$ controls
which degrees appear. For instance, a vector $\times$ vector product ($\ell = 1
\times \ell = 1$) produces $\ell \in \{0, 1, 2\}$: a scalar (the dot product), a
pseudovector (the cross product), and a symmetric traceless tensor.

\paragraph{Tensor Field Networks.}
The Tensor Field Networks of \citet{thomas2018tensorfieldnetworks} provide the
foundational $\SE(3)$-equivariant convolution. A convolution maps an input
tensor field of degree $\ell_{\text{in}}$ to an output of degree
$\ell_{\text{out}}$ through a kernel that depends on the relative position
$\mathbf{r}_{ij} = \mathbf{x}_i - \mathbf{x}_j$, factorized into a learnable
radial part and a fixed angular part given by spherical harmonics:
\begin{equation}
  K^{(\ell_{\text{out}}, \ell_{\text{in}})}_{ij}
  = \sum_{\ell_f} R^{\ell_f}_{\ell_{\text{out}}, \ell_{\text{in}}}(r_{ij})\,
    \mathcal{Y}^{(\ell_f)}(\hat{\mathbf{r}}_{ij}),
  \label{eq:groups:tfn-kernel}
\end{equation}
where $r_{ij} = \norm{\mathbf{r}_{ij}}$, $\hat{\mathbf{r}}_{ij} =
\mathbf{r}_{ij}/r_{ij}$, the $\mathcal{Y}^{(\ell_f)}$ are spherical harmonics,
and the $R^{\ell_f}_{\ell_{\text{out}}, \ell_{\text{in}}}$ are learnable radial
functions (Gaussian shells, Bessel functions, or polynomial envelopes).

\begin{proposition}[Equivariance of the factorized kernel]
\label{prop:groups:tfn-kernel-equivariance}
The kernel \Cref{eq:groups:tfn-kernel} is $\SE(3)$-equivariant: under a rotation
$R \in \SO(3)$,
\begin{equation}
  K^{(\ell_{\text{out}}, \ell_{\text{in}})}_{ij}(R\mathbf{r}_{ij})
  = D^{(\ell_{\text{out}})}(R)\,
    K^{(\ell_{\text{out}}, \ell_{\text{in}})}_{ij}(\mathbf{r}_{ij})\,
    \big(D^{(\ell_{\text{in}})}(R)\big)^{-1}.
\end{equation}
\end{proposition}

\begin{proof}
The radial part is rotation-invariant ($\norm{R\mathbf{r}} = \norm{\mathbf{r}}$),
and the spherical harmonics transform as $Y_\ell^m(R^{-1}\hat{\mathbf{r}}) =
\sum_{m'} D^{(\ell)}_{m'm}(R)\, Y_\ell^{m'}(\hat{\mathbf{r}})$
(\Cref{thm:groups:sh-transformation}). Coupling the harmonics through the
Clebsch--Gordan coefficients translates this transformation into the tensor
transformation law $D^{(\ell_{\text{out}})}(R)\, K\,
(D^{(\ell_{\text{in}})}(R))^{-1}$, as required for equivariance. Translations do
not affect the kernel, which depends only on the relative position.
\end{proof}

The coupling between an input tensor and the angular part is performed through
the Clebsch--Gordan product, with the selection rule $|\ell_{\text{in}} -
\ell_f| \leq \ell_{\text{out}} \leq \ell_{\text{in}} + \ell_f$ determining the
allowed couplings. A complete TFN layer sums these couplings over input degrees,
filter degrees, and neighbors:
\begin{equation}
  T_i^{(\ell), \text{out}}
  = \sum_{\ell'} \sum_{j \in \mathcal{N}(i)} \sum_{\ell_f}
    W^{\ell_f}_{\ell, \ell'}(r_{ij})
    \big(T_j^{(\ell')} \otimes \mathcal{Y}^{(\ell_f)}(\hat{\mathbf{r}}_{ij})
    \big)^{(\ell)}.
\end{equation}

\paragraph{$\SE(3)$-Transformers.}
The $\SE(3)$-Transformers of \citet{fuchs2020se3transformers} combine these
tensor-field convolutions with attention, yielding a more expressive
architecture than EGNNs. The attention weights are computed from invariant
quantities only --- the scalar ($\ell = 0$) components and the pairwise distance:
\begin{equation}
  \alpha_{ij} = \softmax_j\Big(\MLP\big(\mathbf{f}_i^{(0)} \oplus
    \mathbf{f}_j^{(0)} \oplus h(r_{ij})\big)\Big),
  \label{eq:groups:se3-attention}
\end{equation}
where $r_{ij} = \norm{\mathbf{x}_i - \mathbf{x}_j}$, $h(r_{ij})$ is a radial
embedding, and $\oplus$ denotes concatenation.

\begin{proposition}[Invariance of the attention weights]
\label{prop:groups:se3-attention-invariance}
The weights $\alpha_{ij}$ in \Cref{eq:groups:se3-attention} are
$\SE(3)$-invariant.
\end{proposition}

\begin{proof}
The weights depend on positions only through $\mathbf{x}_i - \mathbf{x}_j$ and
its norm. Under $(R, \mathbf{t}) \in \SE(3)$, the difference becomes
$R(\mathbf{x}_i - \mathbf{x}_j)$ and the norm is preserved by orthogonality.
Since $\alpha_{ij}$ is a function of $\norm{\mathbf{x}_i - \mathbf{x}_j}$ and the
invariant scalar features, it is invariant.
\end{proof}

The feature update applies these invariant weights to value messages built by
TFN-style equivariant kernels, $\mathbf{f}_i^{(\ell),\text{new}} = \sum_{j}
\alpha_{ij}\, \mathbf{v}_j^{(\ell)}$, where each value $\mathbf{v}_j^{(\ell)}$
transforms as a degree-$\ell$ tensor.

\begin{theorem}[Equivariance of $\SE(3)$-Transformers]
\label{thm:groups:se3-transformer-equivariance}
The update $\mathbf{f}_i^{(\ell),\text{new}} = \sum_{j} \alpha_{ij}\,
\mathbf{v}_j^{(\ell)}$ with invariant weights \Cref{eq:groups:se3-attention} is
$\SE(3)$-equivariant: $(\mathbf{f}_i^{(\ell),\text{new}})' = D^{(\ell)}(R)\,
\mathbf{f}_i^{(\ell),\text{new}}$.
\end{theorem}

\begin{proof}
The weights $\alpha_{ij}$ are invariant
(\Cref{prop:groups:se3-attention-invariance}), and the value messages, built
from TFN kernels \citep{thomas2018tensorfieldnetworks}, transform as
$(\mathbf{v}_j^{(\ell)})' = D^{(\ell)}(R)\, \mathbf{v}_j^{(\ell)}$. Hence
\begin{align}
  (\mathbf{f}_i^{(\ell),\text{new}})'
  = \sum_{j} \alpha_{ij}\, D^{(\ell)}(R)\, \mathbf{v}_j^{(\ell)}
  = D^{(\ell)}(R) \sum_{j} \alpha_{ij}\, \mathbf{v}_j^{(\ell)}
  = D^{(\ell)}(R)\, \mathbf{f}_i^{(\ell),\text{new}},
\end{align}
using invariance of $\alpha_{ij}$ and linearity of $D^{(\ell)}(R)$. Translations
act only on positions, not on tensor features, completing $\SE(3)$-equivariance.
\end{proof}

The two families occupy complementary points on the expressivity--cost
spectrum, summarized in \Cref{tab:groups:so3-vs-se3}. The compact group $\SO(3)$
admits a finite normalizable Haar measure and pure tensor fields with no
intrinsic invariants, whereas the non-compact $\SE(3)$ has an infinite Haar
measure and natural invariants (relative distances) that allow scalar processing
with standard multilayer perceptrons. EGNNs are cheaper and scale to thousands
of atoms; TFNs and $\SE(3)$-Transformers are more expressive but, with
Clebsch--Gordan products of cost $O(L_{\max}^3)$, are limited to medium systems
with $L_{\max} \leq 4$. The main computational bottleneck of the tensor
approaches is the cost of the Clebsch--Gordan products, which grows with
$L_{\max}$, together with the numerical stability of high-degree spherical
harmonic sums.

\begin{table}[h]
\centering
\caption{Compact $\SO(3)$ versus non-compact $\SE(3)$ equivariant
architectures.}
\label{tab:groups:so3-vs-se3}
\begin{tabular}{@{}lp{4.6cm}p{4.6cm}@{}}
\toprule
\textbf{Property} & \textbf{$\SO(3)$ (compact)} & \textbf{$\SE(3)$ (non-compact)} \\
\midrule
Group structure & Rotations only & Rotations and translations \\
\addlinespace
Compactness & Compact ($\sim \R P^3$) & Non-compact (subgroup $\R^3$ unbounded) \\
\addlinespace
Haar measure & Finite, normalizable & Infinite (non-normalizable on $\R^3$) \\
\addlinespace
Irreducibles & All finite-dimensional, $\dim \rho^{(\ell)} = 2\ell+1$ & Finite ($\SO(3)$ part) and infinite ($\R^3$ part) \\
\addlinespace
Basic invariants & None (transitive on $S^2$) & Distances $\norm{\mathbf{x}_i - \mathbf{x}_j}$ \\
\addlinespace
Architecture & Pure tensor fields & Scalar/vector separation essential \\
\addlinespace
Applications & Spherical data, CMB, climate & Molecules, point clouds, robotics \\
\bottomrule
\end{tabular}
\end{table}

$\Eucl(n)$- and $\SE(3)$-equivariant networks deliver consistent improvements in
data efficiency (often an order-of-magnitude reduction in training
requirements), out-of-distribution generalization, and transferability to unseen
configurations in domains where Euclidean or rigid symmetries are present:
molecular systems, three-dimensional point clouds, and particle analysis.

\section{Permutation groups: learning on sets}
\label{sec:groups:permutations}

The symmetric group $\Sym_n$ is the symmetry of unordered collections: a set has
no intrinsic order, so a network that processes it must be insensitive to
relabeling. Unlike rotations or translations, permutations are purely
combinatorial.

\begin{definition}[Symmetric group]
\label{def:groups:symmetric-group}
The \emph{symmetric group} $\Sym_n$ is the group of all bijections (permutations)
of $\{1, \dots, n\}$ under composition.
\end{definition}

\begin{proposition}[Basic properties of $\Sym_n$]
\label{prop:groups:sn-properties}
$\Sym_n$ has order $n!$, is non-abelian for $n \geq 3$, is generated by the
adjacent transpositions $(i, i+1)$, and admits the Coxeter presentation
\begin{equation}
  \Sym_n = \langle \sigma_1, \dots, \sigma_{n-1} \mid \sigma_i^2 = e,\,
    (\sigma_i \sigma_{i+1})^3 = e,\, (\sigma_i \sigma_j)^2 = e \text{ for }
    |i-j| > 1 \rangle.
\end{equation}
\end{proposition}

\begin{proof}[Partial]
There are $n$ choices for the image of $1$, $n-1$ for $2$, and so on, giving $n!$
permutations. For $n = 3$, $(12)(23) = (123) \neq (132) = (23)(12)$, so
$\Sym_n$ is non-abelian. Generation by transpositions and the presentation are
established by explicit construction; see \citet{dummit2004abstract}.
\end{proof}

For computation, permutations are realized as matrices.

\begin{definition}[Permutation matrix]
\label{def:groups:permutation-matrix}
For $\pi \in \Sym_n$, the permutation matrix $P_\pi \in \R^{n \times n}$ has
$(P_\pi)_{ij} = 1$ if $j = \pi(i)$ and $0$ otherwise.
\end{definition}

\begin{proposition}[Properties of permutation matrices]
\label{prop:groups:permutation-matrix-properties}
Permutation matrices satisfy $P_{\pi_1 \pi_2} = P_{\pi_1} P_{\pi_2}$ (a
homomorphism), $P_\pi^{-1} = P_\pi^\top$ (orthogonality), and $P_\pi \mathbf{x}$
permutes the components of $\mathbf{x}$.
\end{proposition}

\begin{proof}
Writing $(P_\pi)_{ij} = \delta_{i, \pi(j)}$, $(P_{\pi_1} P_{\pi_2})_{ij} =
\sum_k \delta_{i, \pi_1(k)} \delta_{k, \pi_2(j)} = \delta_{i, \pi_1(\pi_2(j))} =
(P_{\pi_1 \pi_2})_{ij}$. Next $(P_\pi^\top)_{ij} = \delta_{j, \pi(i)} =
(P_{\pi^{-1}})_{ij}$, and $P_\pi P_{\pi^{-1}} = P_e = I$ gives $P_\pi^{-1} =
P_\pi^\top$. Finally $(P_\pi \mathbf{x})_i = \sum_j \delta_{i, \pi(j)} x_j =
x_{\pi^{-1}(i)}$.
\end{proof}

\begin{theorem}[Cayley's theorem for $\Sym_n$]
\label{thm:groups:cayley}
Every finite group of order $n$ is isomorphic to a subgroup of $\Sym_n$. In
particular $\Sym_n$ contains a copy of every finite group of order at most $n!$.
\end{theorem}

Cayley's theorem implies that any finite discrete symmetry can be encoded as a
subgroup of permutations, which justifies the general study of
$\Sym_n$-invariant and $\Sym_n$-equivariant architectures.

\paragraph{Invariant and equivariant functions on sets.}
A function $f: \R^{n \times d} \to \R^k$ is \emph{permutation-invariant} if
$f(\pi \cdot X) = f(X)$ for all $\pi \in \Sym_n$, where $\pi \cdot X$ permutes the
rows of $X$; it is \emph{permutation-equivariant} if $f(\pi \cdot X) = \pi \cdot
f(X)$. Basic invariants include the sum $\sum_i \mathbf{x}_i$, the elementwise
maximum, and the mean $\tfrac{1}{n}\sum_i \mathbf{x}_i$.

\paragraph{Representation theory of $\Sym_n$.}
The conjugacy classes of $\Sym_n$ are in bijection with the partitions $\lambda
\vdash n$: two permutations are conjugate if and only if they have the same cycle
structure. The irreducible representations are constructed combinatorially from
\emph{Young diagrams}.

\begin{definition}[Young diagram and standard tableau]
\label{def:groups:young}
A \emph{Young diagram} of shape $\lambda \vdash n$ is an array of $n$ boxes in
left-justified rows, with $\lambda_i$ boxes in row $i$. A \emph{standard Young
tableau} fills the boxes with $1, \dots, n$ so that entries increase along each
row and down each column.
\end{definition}

\begin{theorem}[Irreducible representations of $\Sym_n$]
\label{thm:groups:sn-irreps}
For each partition $\lambda \vdash n$ there is an irreducible representation
$V_\lambda$ of $\Sym_n$ whose dimension equals the number of standard Young
tableaux of shape $\lambda$. Every irreducible representation arises this way,
and $\sum_{\lambda \vdash n} (\dim V_\lambda)^2 = n!$.
\end{theorem}

\begin{proof}[Construction]
For each tableau $T$ of shape $\lambda$, the polytabloid $e_T = a_T b_T$
antisymmetrizes over columns and symmetrizes over rows; $V_\lambda$ is the span
of all polytabloids. The full argument uses the group algebra $\C[\Sym_n]$ and
Schur's lemma; see \citet{sagan2001symmetric} and \citet{fulton1997young}.
\end{proof}

For $\Sym_3$ the partitions and dimensions are $(3)$ (trivial, $\dim 1$), $(2,1)$
(standard, $\dim 2$), and $(1,1,1)$ (sign, $\dim 1$), with $1^2 + 2^2 + 1^2 = 6 =
3!$. Although elegant, explicit convolution over $\Sym_n$ via irreducible
representations is computationally intractable for large $n$, since $|\Sym_n| =
n!$. Practical architectures avoid working with $\Sym_n$ directly and instead
build invariant and equivariant functions elementarily.

\paragraph{DeepSets: universal approximation.}
The central result for permutation-invariant architectures characterizes
completely the form of continuous invariant functions.

\begin{theorem}[DeepSets universal approximation]
\label{thm:groups:deepsets}
Let $f: \mathcal{X}^n \to \mathcal{Y}$ be a continuous permutation-invariant
function with $\mathcal{X} \subset \R^d$ and $\mathcal{Y} \subset \R^k$ compact.
For every $\epsilon > 0$ there exist continuous maps $\phi: \mathcal{X} \to \R^m$
and $\rho: \R^m \to \mathcal{Y}$ (with $m$ large enough) such that
\begin{equation}
  \sup_{X \in \mathcal{X}^n}
    \norm{f(X) - \rho\Big(\sum_{i=1}^n \phi(x_i)\Big)} < \epsilon.
  \label{eq:groups:deepsets}
\end{equation}
Moreover, $\phi$ and $\rho$ can be realized by neural networks
\citep{zaheer2017deepsets}.
\end{theorem}

\begin{proof}
The argument has three steps. \emph{Step 1 (reduction to measures):} identify the
unordered set $\{x_1, \dots, x_n\}$ with the empirical measure $\mu_X =
\tfrac{1}{n}\sum_i \delta_{x_i}$, so $f$ induces a function $\tilde f$ on discrete
probability measures. \emph{Step 2 (moments):} by a Weierstrass-type theorem on
measure spaces, every continuous function on probability measures is approximated
by $\tilde f(\mu) \approx F(\int g_1\, d\mu, \dots, \int g_M\, d\mu)$ for basis
functions $g_1, \dots, g_M$ and a continuous $F$; for the empirical measure,
$\int g_j\, d\mu_X = \tfrac{1}{n}\sum_i g_j(x_i)$. \emph{Step 3 (construction):}
set $\phi(x) = (g_1(x), \dots, g_M(x))$ and $\rho(z_1, \dots, z_M) = F(z_1/n,
\dots, z_M/n)$. Then $\rho(\sum_i \phi(x_i)) = F(\tfrac1n\sum_i g_1(x_i), \dots)
\approx \tilde f(\mu_X) = f(X)$. The universal approximation theorem for neural
networks lets both $\phi$ and $\rho$ be realized by networks.
\end{proof}

\begin{remark}[Latent dimension]
\label{rem:groups:latent-dim}
\citet{zaheer2017deepsets} also prove an exact representation $f(X) =
\rho(\sum_i \phi(x_i))$ for a countable universe, but with possibly
discontinuous $\phi, \rho$. \citet{wagstaff2019limitations} showed that for an
exact representation with continuous functions the latent dimension must satisfy
$m \geq n$. The sum is the unique aggregator guaranteeing universal
approximation: $\max$ and $\text{mean}$ have strictly weaker expressive power.
For example, $\max$ cannot distinguish $\{(1,0),(0,1)\}$ from $\{(1,1)\}$ under
$\phi(x) = x$.
\end{remark}

\paragraph{PointNet.}
PointNet \citep{qi2017pointnet} brings permutation invariance to
three-dimensional point clouds, processing $\mathcal{P} = \{\mathbf{p}_1, \dots,
\mathbf{p}_n\} \subset \R^3$ through a shared per-point encoder, a symmetric
pooling, and a global decoder:
\begin{equation}
  f_{\text{PointNet}}(\mathcal{P})
  = \rho_\psi\Big(\max_{i=1,\dots,n} \phi_\theta(T \mathbf{p}_i)\Big),
  \label{eq:groups:pointnet}
\end{equation}
where $T$ is an input transformation predicted by a T-Net.

\begin{proposition}[Invariance of PointNet]
\label{prop:groups:pointnet-invariance}
The map \Cref{eq:groups:pointnet} is exactly permutation-invariant.
\end{proposition}

\begin{proof}
For $\pi \in \Sym_n$,
\begin{align}
  f_{\text{PointNet}}(\pi \cdot \mathcal{P})
  = \rho_\psi\Big(\max_i \phi_\theta(T \mathbf{p}_{\pi(i)})\Big)
  = \rho_\psi\Big(\max_j \phi_\theta(T \mathbf{p}_j)\Big)
  = f_{\text{PointNet}}(\mathcal{P}),
\end{align}
since the maximum is a symmetric operation.
\end{proof}

Unlike DeepSets, PointNet uses max-pooling rather than summation: although the
maximum has lower theoretical expressive power
(\Cref{rem:groups:latent-dim}), it performs well on point clouds where extremal
geometry is informative. Its limitation --- treating each point independently
until the global pooling, discarding local structure --- motivated PointNet++
\citep{qi2017pointnetplus}, which introduces hierarchical local processing while
preserving permutation invariance.

\paragraph{Set Transformers.}
Global aggregation by sum or max discards information about relations between
elements. Set Transformers introduce attention so that elements interact while
preserving permutation equivariance. The multihead attention block over a set
$X \in \R^{n \times d}$ is
\begin{equation}
  \text{MAB}(X) = \operatorname{Concat}(\text{head}_1, \dots, \text{head}_h) W^O,
  \quad
  \text{head}_i = \Attn(X W_i^Q, X W_i^K, X W_i^V),
\end{equation}
with $\Attn(Q, K, V) = \softmax(QK^\top / \sqrt{d_k}) V$.

\begin{proposition}[Equivariance of attention]
\label{prop:groups:attention-equivariance}
The block $\text{MAB}$ is permutation-equivariant: $\text{MAB}(\pi \cdot X) =
\pi \cdot \text{MAB}(X)$ for all $\pi \in \Sym_n$.
\end{proposition}

\begin{proof}
The product $QK^\top$ and the subsequent multiplication by $V$ permute rows
consistently with the permutation of $X$, while softmax acts independently on
each row, preserving the permutation. The composition is therefore equivariant.
\end{proof}

To control cost, the Induced Set Attention Block uses $m$ learnable inducing
points $I \in \R^{m \times d}$, reducing complexity from $O(n^2)$ to $O(nm)$ with
$m \ll n$:
\begin{equation}
  \text{ISAB}_m(X) = \text{MAB}(X, \text{MAB}(I, X)),
\end{equation}
where $\text{MAB}(A, B)$ denotes attention with queries from $A$ and keys/values
from $B$. Pooling by Multihead Attention uses $k$ learnable seeds $S \in \R^{k
\times d}$ to produce a fixed-size, permutation-invariant output:
$\text{PMA}_k(X) = \text{MAB}(S, \text{rFF}(X))$, with $\text{rFF}$ a row-wise
feedforward. These set-attention mechanisms are the permutation-equivariant
version of the attention introduced in \Cref{ch:grids}.

\begin{example}[DeepSets on a small set]
\label{ex:groups:deepsets-numerical}
Take $\mathcal{X} = \{(2,1),(1,3),(3,2)\}$ with the linear encoder $\phi(\mathbf{x})
= W\mathbf{x} + \mathbf{b}$, $W = \left(\begin{smallmatrix} 1 & -1 \\ 0 & 2 \\ 1 &
1 \end{smallmatrix}\right)$, $\mathbf{b} = (0,1,0)^\top$. Then $\phi((2,1)) =
(1,3,3)^\top$, $\phi((1,3)) = (-2,7,4)^\top$, $\phi((3,2)) = (1,5,5)^\top$, and
the symmetric aggregation gives $\sum_i \phi(\mathbf{x}_i) = (0,15,12)^\top$.
With decoder $\rho(\mathbf{z}) = V\mathbf{z}$, $V = \left(\begin{smallmatrix} 1 & 0
& 0 \\ 0 & 0.1 & 0 \end{smallmatrix}\right)$, the output is $f(\mathcal{X}) =
(0, 1.5)^\top$. Permuting the set leaves the sum --- and hence $f$ ---
unchanged, confirming permutation invariance.
\end{example}

Permutation-invariant and -equivariant architectures have transformed domains
where data lack a natural order: classification, segmentation, and detection on
LiDAR point clouds; molecular property prediction and generation, where atom
labeling is arbitrary; jet identification and anomaly detection in high-energy
physics, where the order of detected particles carries no information; and
abstract reasoning over sets. They are also the conceptual foundation of the
graph neural networks of \Cref{ch:graphs}, where invariance to node relabeling is
the defining symmetry.

\section{Equivariant attention}
\label{sec:groups:equivariant-attention}

The same principle that produces group convolution also produces
\emph{equivariant attention}: building geometric symmetries into the attention
mechanism so that adaptive, content-dependent weights respect the structure of
the data. We adapt the attention of \Cref{ch:grids} to act equivariantly under a
group $\group$.

\begin{theorem}[Construction of equivariant attention]
\label{thm:groups:equivariant-attention}
For a group $\group$ acting on the feature space, an attention mechanism is
$\group$-equivariant if and only if (i) the compatibility functions are
$\group$-invariant and (ii) the query, key, and value transformations commute
with the group action.
\end{theorem}

\begin{proof}[Sketch]
($\Leftarrow$) If the compatibility $a(\mathbf{q}_i, \mathbf{k}_j)$ is invariant,
then $a(\rho(g)\mathbf{q}_i, \rho(g)\mathbf{k}_j) = a(\mathbf{q}_i, \mathbf{k}_j)$,
so the attention weights $\alpha_{ij}$ are invariant; if additionally $W_V$
commutes with $\group$ ($W_V \rho(g) = \rho(g) W_V$), then $\sum_j \alpha_{ij} W_V
\rho(g)\mathbf{x}_j = \rho(g) \sum_j \alpha_{ij} W_V \mathbf{x}_j$, which is
equivariance. ($\Rightarrow$) If the attention is equivariant, the weights must
be invariant under the simultaneous action on queries and keys, forcing the
compatibility functions to be invariant; equivariance of the output with fixed
weights then forces $W_V$ to commute with $\rho(g)$.
\end{proof}

\begin{lemma}[Sufficient condition for self-attention equivariance]
\label{lem:groups:self-attention-equivariance}
Let $\group$ act on $\R^{n \times d}$ by permutations. If the projection
matrices satisfy $\mathbf{W}_Q = \mathbf{W}_K = \mathbf{W}_V = \mathbf{W}$ and
$\mathbf{W}$ commutes with the group action, then self-attention is
$\group$-equivariant.
\end{lemma}

\begin{proof}
For a permutation matrix $P \in \group$ and input $\mathbf{X}$,
\begin{align}
  \text{SelfAttention}(P\mathbf{X})
  &= \softmax\!\Big(\tfrac{(P\mathbf{X}\mathbf{W})(P\mathbf{X}\mathbf{W})^\top}
     {\sqrt{d_k}}\Big)(P\mathbf{X}\mathbf{W}) \\
  &= \softmax\!\Big(\tfrac{P(\mathbf{X}\mathbf{W})(\mathbf{X}\mathbf{W})^\top
     P^\top}{\sqrt{d_k}}\Big)P(\mathbf{X}\mathbf{W})
   = P \cdot \text{SelfAttention}(\mathbf{X}),
\end{align}
where the last step uses commutativity and the equivariance of softmax to
permutations.
\end{proof}

\paragraph{Irreducible attention via Clebsch--Gordan.}
For compact Lie groups, the representation theory of \Cref{sec:groups:tfn}
underlies equivariant attention. Features in each layer decompose by
representation type, $\mathbf{f} = \bigoplus_{\ell=0}^{L_{\max}}
\mathbf{f}^{(\ell)}$ with $\mathbf{f}^{(\ell)} \in \R^{n_\ell \times (2\ell+1)}$,
and the spherical harmonics $Y_\ell^m$ realize the $(2\ell+1)$-dimensional
irreducible representations of $\SO(3)$, transforming under the Wigner matrices
$D^{(\ell)}(R)$ \citep{brocker1985representations}. Combining features of
different types requires the Clebsch--Gordan product
(\Cref{thm:groups:clebsch-gordan}); for instance $D^{(1)} \otimes D^{(1)} =
D^{(0)} \oplus D^{(1)} \oplus D^{(2)}$ yields a scalar (dot product), a
pseudovector (cross product), and a symmetric traceless tensor. An equivariant
attention score combines queries and keys of different types through these
products,
\begin{equation}
  e_{ij} = \sum_{\ell_q, \ell_k, \ell_{\text{out}}}
    \MLP_{\ell_q, \ell_k, \ell_{\text{out}}}
    \big(\mathbf{q}_i^{(\ell_q)} \otimes \mathbf{k}_j^{(\ell_k)}
    \big)^{(\ell_{\text{out}})},
\end{equation}
with the multilayer perceptrons acting independently within each representation
type.

\paragraph{Equivariant transformer architectures.}
Several architectures instantiate this framework. The $\SE(3)$-Transformers of
\citet{fuchs2020se3transformers}, already analyzed in \Cref{sec:groups:tfn},
compute invariant attention weights from scalar features, distances, and the
inner products $\mathbf{f}_{1,i} \cdot \mathbf{f}_{1,j}$ of vector features, then
combine equivariant value messages; the equivariance follows exactly as in
\Cref{thm:groups:se3-transformer-equivariance}. The Equiformer of
\citet{liao2022equiformer} replaces the dot-product score $\mathbf{q}_i^\top
\mathbf{k}_j$ with an equivariant multilayer-perceptron attention $e_{ij} =
\MLP_{\text{attn}}(\mathbf{q}_i \otimes \mathbf{k}_j, \norm{\mathbf{r}_{ij}})$,
where $\otimes$ is the tensor product of irreducible features, enabling richer
interactions than a single scalar score while preserving $\SO(3)$ equivariance.
The LieTransformer of \citet{hutchinson2020lietransformer} provides a general
framework for arbitrary Lie groups by lifting features to the group, applying
attention in the group space, and projecting back,
\begin{equation}
  \text{LieAttn}(\mathbf{x}) = \text{Project}\big(
    \Attn(\text{Lift}(\mathbf{x}))\big),
\end{equation}
which guarantees equivariance through the left regular representation $(L_g
f)(h) = f(g^{-1} h)$. Extending the principle to multi-domain pretraining, the
Equivariant Pretrained Transformer of \citet{jiao2024ept} is a foundation model
unifying small molecules and proteins under a single $\Eucl(3)$-equivariant
framework, exploiting that equivariance is an architectural property, preserved
under transfer to any domain that admits the same group action.

\paragraph{Vector attention on point clouds.}
A complementary line trades exact $\SO(3)$ equivariance for efficiency. The Point
Transformer of \citet{zhao2021point} uses \emph{vector} attention, in which the
attention weight is a vector that modulates feature channels independently:
\begin{equation}
  \mathbf{y}_i = \sum_{j \in \mathcal{N}(i)}
    \rho\big(\gamma(\phi(\mathbf{x}_i) - \psi(\mathbf{x}_j) + \delta)\big)
    \big(\alpha(\mathbf{x}_j) + \delta\big),
\end{equation}
where $\phi, \psi, \alpha$ are linear maps, $\gamma, \rho$ are multilayer
perceptrons, and $\delta = \MLP(\mathbf{p}_i - \mathbf{p}_j)$ encodes the relative
position. Because $\delta$ depends only on the relative position $\mathbf{p}_i -
\mathbf{p}_j$, the attention weights are translation-invariant, connecting
directly to the invariance condition of
\Cref{thm:groups:equivariant-attention}. Avoiding the Clebsch--Gordan products
of cost $O(\ell_{\max}^3)$, the Point Transformer attains complexity
$O(|\mathcal{N}(i)| \cdot d)$ per point, at the price of giving up exact $\SO(3)$
equivariance of the output features.

\section{A unified view}
\label{sec:groups:unified}

The architectures of this chapter realize a single principle in increasingly
general settings. The choice of symmetry group determines both the
representational capacity and the computational cost, summarized in
\Cref{tab:groups:comparison}.

\begin{table}[h]
\centering
\small
\caption{Symmetry groups, complexity, and behavior of equivariant
architectures. $N$: spatial resolution; $L$: maximum spherical degree; $F$:
maximum harmonic frequency.}
\label{tab:groups:comparison}
\begin{tabular}{@{}lcccl@{}}
\toprule
\textbf{Group} & \textbf{Order} & \textbf{Type} & \textbf{Complexity} & \textbf{Domain} \\
\midrule
Standard CNN & --- & Translational & $O(N^2)$ & Grids \\
$C_4$ & $4$ & Finite & $O(4 N^2)$ & Discrete rotations \\
$D_4$ & $8$ & Finite & $O(8 N^2)$ & Rotations + reflections \\
$\SO(2)$ steerable & $\infty$ & Continuous & $O(F^2 N^2)$ & Continuous planar rotations \\
$\SO(3)$ spherical & $\infty$ & Continuous & $O(L^3 N^2)$ & Spherical data \\
$\SE(3)$ & $\infty$ & Non-compact & $O(L^4 N^2)$ & Point clouds, molecules \\
$\Sym_n$ & $n!$ & Finite & $O(n^2)$ & Sets, graphs \\
\bottomrule
\end{tabular}
\end{table}

Three guidelines follow from the trade-off between equivariance fidelity and
computational cost. For real-time settings, discrete group convolutional
networks ($C_4$, $D_4$) give the best balance of accuracy gain and overhead. For
high accuracy, steerable and $\SE(3)$ architectures justify their greater cost
with superior precision. When the symmetries are unknown, one should start with
simple group convolutional networks to confirm their presence before adopting
heavier machinery, following the iterative path: standard convolutional network
$\to$ discrete group convolutional network $\to$ continuous extensions.

The single thread running through the chapter is that equivariance turns a
symmetry of the data into a constraint on the architecture, and that constraint
yields parameter efficiency, robust generalization, and a constructive recipe
for whole networks. The group convolutions developed here generalize classical
convolution from translations to arbitrary symmetry groups, but they still assume
\emph{global}, exact algebraic symmetries. The next step in the hierarchy of
geometric deep learning is to handle data on general geometric spaces, where the
symmetries are local or only approximate. Discarding the ambient grid entirely
and keeping only relational structure leads to the graph neural networks of
\Cref{ch:graphs}; endowing the domain with intrinsic distance and curvature leads
to the manifolds and geodesics of \Cref{ch:geodesics}; and recognizing that a
curved domain admits no canonical local frame leads to the gauge-equivariant
networks of \Cref{ch:gauges}.
